%==============================================================================
% Literature Review (Accessible Version)
% Hybrid Model-Based and Data-Driven Score Fusion
% ECE M237: Dynamic Programming, Spring 2026
% Author: Zihao Sun (UCLA ECE)
%
% Engineering-friendly version assuming multivariable calculus, linear
% algebra, and basic probability. No measure theory.
% Rigorous version: literature_review.tex
%==============================================================================
\documentclass[11pt,letterpaper]{article}

\usepackage[margin=1in]{geometry}
\usepackage[T1]{fontenc}
\usepackage{lmodern}
\usepackage{microtype}
\usepackage{parskip}
\usepackage{amsmath,amssymb,amsthm,amsfonts}
\usepackage{mathtools}
\usepackage{mathrsfs}
\usepackage{bm}
\usepackage{dsfont}
\usepackage{graphicx}
\usepackage{booktabs}
\usepackage{array}
\usepackage{enumitem}
\usepackage[numbers,sort&compress]{natbib}
\usepackage[dvipsnames,table]{xcolor}
\usepackage{algorithm}
\usepackage[noend]{algpseudocode}
\usepackage[framemethod=TikZ]{mdframed}
\usepackage{hyperref}
\usepackage[capitalize,noabbrev]{cleveref}

% ── Skippable advanced-math environment ──────────────────────────────
\newmdenv[
  backgroundcolor=gray!6,
  linecolor=gray!40,
  linewidth=0.8pt,
  roundcorner=2pt,
  innertopmargin=10pt,
  innerbottommargin=10pt,
  skipabove=10pt,
  skipbelow=10pt,
  frametitle={\small\sffamily\bfseries
    $\blacktriangleright$ Advanced Mathematics
    \normalfont\small\itshape\ --- can be skipped on first reading},
  frametitlebackgroundcolor=gray!18,
  frametitleaboveskip=4pt,
  frametitlebelowskip=4pt,
]{advancedmath}
\hypersetup{colorlinks=true,citecolor=blue!60!black,linkcolor=blue!60!black,
            urlcolor=blue!60!black}

\theoremstyle{plain}
\newtheorem{theorem}{Theorem}[section]
\newtheorem{proposition}[theorem]{Proposition}
\newtheorem{lemma}[theorem]{Lemma}
\newtheorem{corollary}[theorem]{Corollary}
\theoremstyle{definition}
\newtheorem{definition}[theorem]{Definition}
\newtheorem{assumption}[theorem]{Assumption}
\newtheorem{example}[theorem]{Example}
\theoremstyle{remark}
\newtheorem{remark}[theorem]{Remark}

\newcommand{\R}{\mathbb{R}}
\newcommand{\N}{\mathbb{N}}
\newcommand{\E}{\mathbb{E}}
\newcommand{\Prob}{\mathbb{P}}
\newcommand{\Var}{\operatorname{Var}}
\newcommand{\ind}{\mathbf{1}}
\newcommand{\dd}{\mathrm{d}}
\newcommand{\KL}{D_{\mathrm{KL}}}
\newcommand{\calB}{\mathcal{B}}
\newcommand{\calF}{\mathcal{F}}
\newcommand{\calL}{\mathcal{L}}
\newcommand{\calN}{\mathcal{N}}
\newcommand{\calP}{\mathcal{P}}
\newcommand{\calT}{\mathcal{T}}
\newcommand{\calU}{\mathcal{U}}
\newcommand{\calX}{\mathcal{X}}
\newcommand{\calH}{\mathcal{H}}
\newcommand{\tr}{\operatorname{tr}}
\newcommand{\grad}{\nabla}
\newcommand{\divg}{\nabla\cdot}
\newcommand{\inner}[2]{\langle #1,\,#2\rangle}
\newcommand{\norm}[1]{\lVert #1\rVert}
\DeclareMathOperator*{\argmin}{arg\,min}
\DeclareMathOperator*{\argmax}{arg\,max}
\newcommand{\smodel}{s_{\mathrm{mod}}}

\title{\bfseries
Literature Review: Score-Based Generative Models, \\
Flow Matching, and Dynamic Programming for Robot Trajectory Generation}
\author{Zihao Sun\\[2pt]
  \small UCLA ECE --- ECE M237: Dynamic Programming (Spring 2026)}
\date{\today}

\begin{document}
\maketitle

\begin{abstract}
We survey the mathematical foundations of three converging threads in
robot trajectory generation: (i)~dynamic programming and optimal control,
(ii)~score-based diffusion and flow matching, and (iii)~stochastic
optimal control via path-integral methods.  The main text is written for
engineers with solid calculus, linear algebra, and basic probability.
Sections containing rigorous measure-theoretic derivations and full
proofs are enclosed in shaded
\textbf{``Advanced Mathematics''} boxes and can be safely
skipped on a first reading without losing the narrative thread.
A self-contained \textbf{Appendix} reproduces the full
measure-theoretic treatment---It\^o calculus, Girsanov's theorem,
Feynman--Kac, Bellman contraction, and complete proofs of all bridge
theorems---for reference in future course projects.
\end{abstract}

\tableofcontents
\bigskip

%==============================================================================
\section{Introduction}\label{sec:intro}
%==============================================================================

\subsection{The robot trajectory generation problem}

A robot must plan a trajectory $\tau=(x_0,u_0,x_1,u_1,\ldots,x_H)$
consisting of states~$x_k$ and controls~$u_k$ over a planning
horizon~$H$.  A good trajectory must satisfy several competing
objectives simultaneously:
\begin{itemize}[nosep]
  \item \textbf{Physical feasibility:} the trajectory must obey the
    robot's dynamics $x_{k+1}=f(x_k,u_k)$ (joint limits, torque limits,
    contact constraints).
  \item \textbf{Cost minimization:} a task cost $J(\tau)$ encodes goals
    such as reaching a target, minimizing energy, or following a
    reference path.
  \item \textbf{Safety:} obstacles and forbidden regions must be avoided
    at every time step; formally, a safety set
    $\{x:h(x)\ge 0\}$ must remain forward-invariant.
  \item \textbf{Naturalness and diversity:} in human-interactive
    settings, trajectories should reflect human preferences and capture
    the multimodality of acceptable behaviors (e.g., reaching around an
    obstacle from the left or right).
\end{itemize}
Classical model-based methods (trajectory optimization, MPC) excel at
feasibility and cost minimization but produce unimodal solutions and
require differentiable dynamics.  Purely data-driven methods
(behavioral cloning, offline RL) capture diversity but ignore physics
at test time, producing trajectories that may violate dynamics or
collide with novel obstacles.  Our project bridges these two worlds.

\subsection{Motivation: two paradigms and their limitations}

\textbf{Data-driven score models} learn a trajectory distribution
$p^{\mathrm{data}}$ from expert demonstrations via score-based SDEs
\citep{song2021score} or flow-matching ODEs
\citep{lipman2023flow,chi2023diffusionpolicy,janner2022}.
Diffusion-based planners such as Diffuser~\citep{janner2022} and
flow-matching methods such as FlowMP~\citep{flowmp} learn expressive
multimodal priors from demonstrations, but provide no feasibility
guarantees out of distribution.  Kernel-based Behavioral Score
Diffusion~\citep{bsd} offers a training-free data-driven variant.
These methods capture multimodal behavior but ignore robot dynamics and
task cost at test time.

\textbf{Model-based score methods.}  Model-Based Diffusion
(MBD)~\citep{mbd} derives the score analytically from known dynamics
$x_{k+1}=f(x_k,u_k)$ and cost $J(\tau)$ via Monte Carlo rollouts,
guaranteeing feasibility without any learned component.  Its authors
note that data ``can be naturally integrated'' but leave this
unformalized---the gap our project addresses.  MBD requires no training
data but cannot express human preference distributions and suffers
high variance when the trajectory dimension is large or the temperature
$\lambda$ is small.

\subsection{Three bodies of literature}

Our project sits at the intersection of three research threads:
\begin{enumerate}[nosep]
  \item \textbf{Model-based score computation:}
    MBD~\citep{mbd} and SOCM~\citep{domingoenrich2024socm} compute or
    learn scores from dynamics and cost, providing feasibility
    guarantees.
  \item \textbf{Generative trajectory planning:}
    Diffuser~\citep{janner2022}, Diffusion
    Policy~\citep{chi2023diffusionpolicy}, FlowMP~\citep{flowmp},
    and BSD~\citep{bsd} learn expressive priors from demonstrations
    via diffusion or flow matching.
  \item \textbf{Safe generative control:}
    SafeFlow~\citep{safeflow}, SafeDiffuser~\citep{xiao2025safediffuser},
    $S^2$Diff~\citep{s2diff}, and related methods enforce safety via
    control barrier functions (CBFs) or Lyapunov conditions.
\end{enumerate}
Further connections include classifier-free guidance as policy
improvement~\citep{cfgrl} and optimal-control-guided flow
matching~\citep{ocflow}.

\subsection{Our approach: adaptive product-of-experts fusion}

We fuse data-driven and model-based scores via a product-of-experts:
\begin{equation}
  \grad\log p^{\mathrm{target}}_t(\tau)
  \;=\; s_\theta(\tau,t)
  + \lambda(x,t)\,\smodel(\tau,t),
  \label{eq:poe-intro}
\end{equation}
where the state- and time-dependent weight $\lambda(x,t)$ is designed
to \emph{trust the data prior} in regions well covered by the training
distribution, and to \emph{rely on the physics model} in
out-of-distribution states.  This adaptive weighting mechanism---and
the empirical study of the tradeoff it governs---is the primary novel
contribution of our work.

We further conjecture that in the linear-quadratic (LQ) special case,
the model-based correction analytically reduces to the LQR gain
$-\lambda^{-1}P_k x_k$ governed by the Discrete Algebraic Riccati
Equation (DARE), yielding a formal bridge between generative trajectory
modeling and classical optimal control.

\subsection{Conventions}
We work with smooth probability densities. Gradients $\grad$ are
w.r.t.\ the state variable. Technical regularity conditions and full
proofs are in the shaded \textbf{Advanced Mathematics} boxes
throughout the document; these can be safely skipped on first reading.

%==============================================================================
\section{Stochastic Calculus Essentials}\label{sec:prelim}
%==============================================================================

The results in this section provide the mathematical toolkit
underlying score-based models and stochastic optimal control. The key
takeaways are: (i)~It\^o's formula adds a second-order correction to
the chain rule; (ii)~the Fokker--Planck equation describes how
densities evolve under an SDE; (iii)~Feynman--Kac converts PDEs to
Monte Carlo problems; (iv)~Girsanov's theorem links KL divergence to
control energy; and (v)~reversing a diffusion requires the score
$\grad\log p_t$.

\begin{advancedmath}

\subsection{Brownian motion}

\begin{definition}
A $d$-dimensional Brownian motion $W_t$ starts at $W_0=0$ with:
(i) independent increments: $W_t-W_s$ is independent of the past;
(ii) Gaussian increments: $W_t-W_s\sim\calN(0,(t-s)I_d)$;
(iii) continuous paths.
\end{definition}

\textbf{Key property:} $(\dd W_t)^2=\dd t$ (quadratic variation). This
is why stochastic calculus has an extra correction term.

\subsection{It\^o's formula: the stochastic chain rule}

\begin{theorem}[It\^o's formula]\label{thm:ito}
If $X_t$ satisfies $\dd X_t=\mu_t\dd t+\Sigma_t\dd W_t$ and $\varphi$
is smooth, then:
\begin{equation}
  \dd\varphi(t,X_t)
  = \Bigl[\frac{\partial\varphi}{\partial t}
    + \mu_t^\top\grad\varphi
    + \tfrac{1}{2}\tr(\Sigma_t\Sigma_t^\top\nabla^2\varphi)\Bigr]\dd t
  + (\grad\varphi)^\top\Sigma_t\,\dd W_t.
  \label{eq:ito-formula}
\end{equation}
\end{theorem}

The $\frac{1}{2}\tr(\Sigma\Sigma^\top\nabla^2\varphi)$ is the ``It\^o
correction''---it arises because $(\dd W)^2=\dd t$ makes a second-order
Taylor term survive.

\subsection{SDEs: existence and uniqueness}

\begin{definition}
An SDE describes a random process via
$\dd X_t = b(t,X_t)\dd t + \sigma(t,X_t)\dd W_t$, $X_0\sim p_0$.
\end{definition}

\begin{theorem}
If $b,\sigma$ are Lipschitz in $x$ with at most linear growth, then the
SDE has a unique solution with continuous paths and finite moments.
\end{theorem}

This is the stochastic Picard--Lindel\"of theorem.

\subsection{Fokker--Planck: how densities evolve}

\begin{theorem}[Fokker--Planck]\label{thm:fp}
If $X_t$ solves the SDE and has density $p_t(x)$, then
\begin{equation}
  \frac{\partial p_t}{\partial t}
  = -\divg(b\,p_t) + \frac{1}{2}\sum_{i,j}\frac{\partial^2}{\partial x_i\partial x_j}
    [(\sigma\sigma^\top)_{ij}\,p_t].
  \label{eq:fp}
\end{equation}
\end{theorem}

The first term is transport (drift pushes probability); the second is
diffusion (noise spreads probability, like heat).

\subsection{Feynman--Kac: PDEs as expectations}

\begin{theorem}[Feynman--Kac]\label{thm:fk}
The PDE $\partial_t u+b^\top\grad u+\frac{1}{2}\tr(\sigma\sigma^\top\nabla^2u)-Vu=0$,
$u(T,x)=f(x)$, has solution
\begin{equation}
  u(t,x) = \E\!\left[f(X_T)\exp\!\Bigl(-\int_t^T V(X_s)\dd s\Bigr)\;\Big|\;X_t=x\right].
  \label{eq:fk}
\end{equation}
\end{theorem}

This converts PDEs into Monte Carlo sampling problems---the key link
between optimal control and sampling.

\subsection{Girsanov: changing drift by reweighting}

\begin{theorem}[Girsanov, informal]\label{thm:girsanov}
Reweighting paths of $\dd X_t=b\dd t+\sigma\dd W_t$ by
$Z_T=\exp(\int_0^T\theta_s^\top\dd W_s-\frac{1}{2}\int_0^T\norm{\theta_s}^2\dd s)$
makes $X_t$ behave as if driven by $\dd X_t=(b+\sigma\theta_t)\dd t+\sigma\dd\tilde W_t$.
\end{theorem}

\textbf{Key consequence:} The KL divergence between a controlled SDE
$\dd X=(b+\sigma u)\dd t+\sigma\dd W$ and the uncontrolled $\dd X=b\dd t+\sigma\dd W$ is
\begin{equation}
  \KL(\text{controlled}\|\text{uncontrolled})=\tfrac{1}{2}\E\!\int_0^T\!\norm{u_t}^2\dd t.
  \label{eq:kl-path}
\end{equation}
KL divergence \emph{equals} control energy. This is the foundation of
path-integral control (\Cref{sec:soc}).

\subsection{Time reversal}

\begin{theorem}[Anderson]\label{thm:time-rev}
The time-reversed process $\bar X_s=X_{T-s}$ of
$\dd X_t=b\dd t+\sigma\dd W_t$ satisfies:
\begin{equation}
  \dd\bar X_s=[-b(T\!-\!s,\bar X_s)+\sigma\sigma^\top\!\grad\log p_{T-s}(\bar X_s)]\dd s
  +\sigma(T\!-\!s)\dd\bar W_s.
  \label{eq:time-rev}
\end{equation}
\end{theorem}

To reverse a diffusion, flip the drift and add a score correction
$\sigma\sigma^\top\grad\log p_t$. This is the mathematical basis of
\emph{all} score-based generative models.

\subsection{Relative entropy and Donsker--Varadhan}

The \emph{KL divergence} (relative entropy) between two distributions is
$\KL(\mu\|\nu)=\int\log(\dd\mu/\dd\nu)\dd\mu$.
The \emph{Donsker--Varadhan variational formula} states:
\[
  -\log\E^\nu[e^{-F}]=\inf_{\mu\ll\nu}\{\E^\mu[F]+\KL(\mu\|\nu)\},
\]
with the infimum attained at $\dd\mu^*/\dd\nu\propto e^{-F}$.
Together with \eqref{eq:kl-path}, this yields the KL-control identity
used throughout stochastic optimal control (\Cref{sec:soc}).

\end{advancedmath}

%==============================================================================
\section{Dynamic Programming and Optimal Control}\label{sec:dp}
%==============================================================================

\subsection{Bellman equation}

For an MDP with cost $g$, transitions $P$, discount $\gamma$:
\begin{equation}
  V^*(x)=\min_u\{g(x,u)+\gamma\E[V^*(x')|x,u]\}.
  \label{eq:bellman}
\end{equation}
Value iteration converges at rate $\gamma^k$ because the Bellman operator
is a $\gamma$-contraction (Banach fixed-point theorem).

\subsection{LQR: closed-form DP}\label{sec:lqr}

\textbf{Setup.} $x_{k+1}=Ax_k+Bu_k+w_k$, $w_k\sim\calN(0,\Sigma_w)$,
cost $J=\E[\sum_{k=0}^{H-1}(x_k^\top Qx_k+u_k^\top Ru_k)+x_H^\top Q_Hx_H]$.

\begin{theorem}[Riccati]\label{thm:lqr}
$V_k^*(x)=x^\top P_kx+c_k$ with
\begin{equation}
  P_k=Q+A^\top P_{k+1}A-A^\top P_{k+1}B(R+B^\top P_{k+1}B)^{-1}B^\top P_{k+1}A,
  \quad P_H=Q_H.
  \label{eq:dare}
\end{equation}
Optimal control: $u_k^*=-K_kx_k$,
$K_k=(R+B^\top P_{k+1}B)^{-1}B^\top P_{k+1}A$.
\end{theorem}

\begin{advancedmath}
\begin{proof}[Derivation]
Backward induction. Assume $V_{k+1}^*=x^\top P_{k+1}x+c_{k+1}$.
Substituting into the Bellman equation gives a quadratic in $u$:
$u^\top(R+B^\top P_{k+1}B)u+2u^\top B^\top P_{k+1}Ax+\dots$ Setting
$\partial/\partial u=0$ yields $K_k$; substituting back gives $P_k$.
$P_k\succeq 0$ follows from a Schur-complement identity.
If $(A,B)$ is stabilizable and $(A,Q^{1/2})$ detectable,
$P_\infty=\lim P_k$ solves the DARE uniquely in $\succeq 0$, and
$A-BK_\infty$ is Schur stable.
\end{proof}

\begin{remark}[Value as log-density]\label{rem:logdensity}
Define $\nu_k(x)\propto\exp(-V_k^*(x)/\lambda)$. This is Gaussian with
$\grad\log\nu_k(x)=-\lambda^{-1}P_kx$. The score of the cost-shaped
Boltzmann distribution is the \emph{negated Riccati gain}. This is the
origin of \Cref{thm:lq-score}.
\end{remark}
\end{advancedmath}

\begin{advancedmath}
\subsection{HJB equation}

For continuous-time $\dd x=f(x,u)\dd t+\sigma(x)\dd W_t$, cost
$J=\E[\int g\dd t+h(x_T)]$:
\begin{equation}
  -\partial_tV=\min_u\{g(x,u)+(\grad V)^\top f(x,u)
    +\tfrac{1}{2}\tr(\sigma\sigma^\top\nabla^2V)\},\quad V(T,x)=h(x).
  \label{eq:hjb}
\end{equation}
Under bounded continuous coefficients and uniform parabolicity
($\sigma\sigma^\top\succeq\epsilon I$), $V$ is the unique bounded
\emph{viscosity solution} of the HJB.

\subsection{Pontryagin Maximum Principle}\label{sec:pmp}

For deterministic $\dot x=f(x,u)$, cost $J=\int g\dd t+h(x_T)$, the
Hamiltonian $\mathcal{H}=g+\lambda^\top f$ yields:
\begin{equation}
  \dot x=\partial_\lambda\mathcal{H},\quad
  \dot\lambda=-\partial_x\mathcal{H},\quad
  u^*=\argmin_u\mathcal{H},\quad
  \lambda_T=\grad h(x_T).
  \label{eq:pmp}
\end{equation}
OC-Flow \citep{ocflow} applies PMP to guide flow matching:
$u^*=-\lambda_t$ with adjoint
$\dot\lambda=-(\partial_xv_\theta)^\top\lambda$.
\end{advancedmath}

%==============================================================================
\section{Score-Based Generative Modeling}\label{sec:score}
%==============================================================================

The core idea behind score-based generative models is remarkably
simple: if we know how to gradually \emph{destroy} structure in data
by adding noise, we can learn to \emph{reverse} that process to
generate new data from pure noise.  Concretely, given a dataset of
trajectories (or images, or any high-dimensional objects), we define a
forward process that progressively corrupts each sample with Gaussian
noise until it becomes indistinguishable from random noise.  Then we
train a neural network to predict the ``score''---the gradient of the
log-density $\grad\log p_t(x)$---at every noise level.  Once trained,
this score network lets us simulate the reverse process, starting from
random noise and iteratively denoising to produce a clean sample.

\textbf{DDPM}~\citep{ho2020ddpm} introduced this framework in
discrete time: a $T$-step forward chain adds noise
$x_t=\sqrt{\bar\alpha_t}\,x_0+\sqrt{1-\bar\alpha_t}\,\epsilon$ (with
$\epsilon\sim\calN(0,I)$), and a learned denoiser $\epsilon_\theta$
predicts the noise to reverse each step.  The training loss is simply
$\E\norm{\epsilon-\epsilon_\theta(x_t,t)}^2$.  The connection to
scores is $\grad\log p_t=-\epsilon/\sqrt{1-\bar\alpha_t}$, so the
noise predictor \emph{is} a scaled score estimator.

\textbf{Score SDE}~\citep{song2021score} unifies DDPM and related
methods under a continuous-time SDE framework, enabling exact
likelihood computation and principled guidance.  We summarize the key
components below.

\begin{algorithm}[ht]
\caption{DDPM Training}\label{alg:ddpm-train}
\begin{algorithmic}[1]
\Require Dataset $\mathcal{D}$, noise schedule $\{\bar\alpha_t\}_{t=1}^T$, network $\epsilon_\theta$
\Repeat
  \State Sample $x_0\sim\mathcal{D}$, $t\sim\mathrm{Uniform}(\{1,\ldots,T\})$, $\epsilon\sim\calN(0,I)$
  \State $x_t\gets\sqrt{\bar\alpha_t}\,x_0+\sqrt{1-\bar\alpha_t}\,\epsilon$ \Comment{Add noise}
  \State Take gradient step on $\norm{\epsilon-\epsilon_\theta(x_t,t)}^2$
\Until{converged}
\end{algorithmic}
\end{algorithm}

\begin{algorithm}[ht]
\caption{DDPM Sampling (Reverse Process)}\label{alg:ddpm-sample}
\begin{algorithmic}[1]
\Require Trained $\epsilon_\theta$, noise schedule $\{\alpha_t,\bar\alpha_t,\sigma_t\}$
\State $x_T\sim\calN(0,I)$
\For{$t=T,T-1,\ldots,1$}
  \State $z\sim\calN(0,I)$ if $t>1$, else $z=0$
  \State $\mu_t\gets\frac{1}{\sqrt{\alpha_t}}\bigl(x_t-\frac{1-\alpha_t}{\sqrt{1-\bar\alpha_t}}\epsilon_\theta(x_t,t)\bigr)$
  \State $x_{t-1}\gets\mu_t+\sigma_tz$ \Comment{Denoise one step}
\EndFor
\State \Return $x_0$
\end{algorithmic}
\end{algorithm}

\begin{advancedmath}
\subsection{Forward SDE: adding noise}

Start with data $p_0$. Add noise:
$\dd X_t=f(t)X_t\dd t+g(t)\dd W_t$. After time $T$,
$p_T\approx\calN(0,\sigma_T^2I)$.

\begin{definition}[Score]
$s_p(x):=\grad\log p(x)=\grad p(x)/p(x)$. It points toward higher
probability. $\E_p[s_p]=0$.
\end{definition}

\subsection{Reverse SDE: generating by denoising}

\begin{theorem}[Anderson \citep{anderson1982reverse}]
The reversed process satisfies
$\dd\bar X_s=[-f\bar X_s+g^2\grad\log p_{T-s}(\bar X_s)]\dd s+g\dd\bar W_s$.
\end{theorem}

\begin{proof}[Proof sketch]
Let $q_s(x):=p_{T-s}(x)$. The Fokker--Planck equation for $p_t$
becomes (with a sign flip from $t=T-s$):
$\partial_s q_s=\divg[(b-\sigma\sigma^\top\grad\log q_s)q_s]
+\tfrac12\divg[\sigma\sigma^\top\grad q_s]$,
which is the Fokker--Planck equation of the reverse SDE.
Weak uniqueness of SDEs determines the law from its marginal PDE.
\end{proof}

If we know the score $\grad\log p_t$ at every noise level, we can
start from noise and simulate backward to get exact data samples.

\subsection{Denoising score matching}

\begin{theorem}[Vincent \citep{vincent2011connection}]\label{thm:dsm}
The intractable loss $\E_{p_t}\norm{s_\theta-\grad\log p_t}^2$ equals
(up to a constant) the computable denoising loss:
\begin{equation}
  \mathcal{L}_{\mathrm{DSM}}
  =\E_{t,x_0,x_t}\norm{s_\theta(x_t,t)-\grad_{x_t}\log p_{t|0}(x_t|x_0)}^2.
  \label{eq:dsm}
\end{equation}
\end{theorem}

\begin{proof}
Expand the square in both sides. The cross terms are equal because
$\grad p_t(x_t)=\int p_0(x_0)p_{t|0}(x_t|x_0)\grad_{x_t}\log p_{t|0}\,\dd x_0$,
so $-\int s_\theta\cdot\grad p_t=-\E_{(X_0,X_t)}[s_\theta^\top\grad\log p_{t|0}]$.
The squared-norm terms that do not involve $\theta$ contribute only a
$\theta$-independent constant.
\end{proof}

For Gaussian noise $p_{t|0}=\calN(\alpha_tx_0,\sigma_t^2I)$, the target
is $\grad\log p_{t|0}=-(x_t-\alpha_tx_0)/\sigma_t^2$: just point back
toward the clean data. This is DDPM's $\epsilon$-prediction loss
\citep{ho2020ddpm}.

\subsection{Probability-flow ODE}

\begin{theorem}[\citet{song2021score}]
The ODE $\dd x_t=[f(t)x_t-\frac{1}{2}g(t)^2\grad\log p_t(x_t)]\dd t$
produces the same marginals $p_t$ as the SDE, without randomness.
\end{theorem}

\begin{proof}
Fokker--Planck gives
$\partial_tp_t=-\divg([b-\tfrac12\sigma\sigma^\top\grad\log p_t]p_t)$,
which is the continuity equation of this ODE. Uniqueness of
continuity-equation solutions with Lipschitz vector fields gives the
result.
\end{proof}

This enables exact likelihoods and bridges diffusion to flow matching.
\end{advancedmath}

%==============================================================================
\section{Flow Matching}\label{sec:fm}
%==============================================================================

Flow matching~\citep{lipman2023flow} is a recent alternative to
diffusion that has rapidly become the preferred generative backbone
for trajectory planning.  While diffusion models add and remove noise
via stochastic processes (SDEs), flow matching learns a
\emph{deterministic} velocity field that transports noise to data
along smooth paths.  This seemingly small difference has major
practical consequences:

\begin{itemize}[nosep]
  \item \textbf{Faster sampling:} flow matching generates samples by
    solving an ODE, which requires only 10--20 neural-network
    evaluations with an adaptive-step solver, compared to 100+ denoising
    steps for DDPM.
  \item \textbf{Simulation-free training:} the training loss is a simple
    regression that does not require simulating the forward or reverse
    process during training.
  \item \textbf{Straight paths:} the default linear interpolant
    $x_t=(1-t)x_0+tx_1$ produces nearly straight transport paths, which
    are easier for the network to learn and faster for ODE solvers to
    integrate.
  \item \textbf{Natural connection to optimal transport:}
    straight-line interpolation can be shown to approximate the optimal
    transport map between the noise and data distributions
    \citep{tong2023ot}, minimizing the total displacement.
\end{itemize}

\subsection{Idea: transport noise to data via a velocity field}

The goal is to learn a velocity field $v_\theta(x,t)$ such that the
ODE $\dot x_t=v_\theta(x_t,t)$ transports the noise distribution
$p_0=\calN(0,I)$ to the data distribution $p_1=p^{\mathrm{data}}$
over the time interval $t\in[0,1]$.  The density along this flow
satisfies the \emph{continuity equation}:
$\partial_tp_t+\divg(p_tv_t)=0$.

Intuitively, think of $v_t$ as a wind field that blows particles of
probability from where they start (noise) to where they should end
(data).  At time $t=0$, particles are spread out as a Gaussian cloud;
at time $t=1$, they have been transported to match the data
distribution.

\subsection{Conditional flow matching (CFM)}

The marginal flow-matching loss
$\E_{t,x\sim p_t}\norm{v_\theta(x,t)-v_t(x)}^2$ is intractable
because computing the true velocity $v_t$ requires knowledge of the
full marginal density $p_t$.  The breakthrough of conditional flow
matching (CFM) is to bypass this by training on \emph{per-sample}
conditional paths.

\begin{theorem}[\citet{lipman2023flow}]\label{thm:cfm}
For straight-line paths $x_t=(1-t)x_0+tx_1$ ($x_0\sim\calN(0,I)$,
$x_1\sim p_{\mathrm{data}}$) with conditional velocity $v_t=x_1-x_0$,
the computable loss
\begin{equation}
  \mathcal{L}_{\mathrm{CFM}}
  =\E_{t,x_0,x_1}\norm{v_\theta(x_t,t)-(x_1-x_0)}^2
  \label{eq:cfm}
\end{equation}
has the same gradients as the intractable marginal loss.
\end{theorem}

This is the flow-matching analog of denoising score matching: we
cannot compute the true marginal target, but training against the
conditional target (which is trivially computable---just the difference
$x_1-x_0$) gives the same optimal network.

\begin{advancedmath}
\begin{proof}[Proof of \Cref{thm:cfm}]
Define the marginal $p_t(x)=\int p_t(x|x_1)q(x_1)\dd x_1$ and
the marginal velocity
$v_t(x)=\E[v_t(x|X_1)\mid X_t=x]=\int v_t(x|x_1)\frac{p_t(x|x_1)q(x_1)}{p_t(x)}\dd x_1$.
Differentiate $p_t$ in $t$ and apply the conditional continuity
equation: $\partial_tp_t(x)=-\int\divg_x[p_t(x|x_1)v_t(x|x_1)]q(x_1)\dd x_1$.
Swapping divergence and integral gives $\partial_tp_t=-\divg(p_tv_t)$,
confirming that the marginal velocity generates $(p_t)$.
For the loss identity, expand both squared norms; the LHS cross term
$-2\E_{t,p_t}\langle v_\theta,v_t\rangle$ equals the RHS cross term by
substituting \eqref{eq:cfm} and using Fubini.
\end{proof}

\textbf{Optimal transport geometry.}
The straight-line interpolant $x_t=(1-t)x_0+tx_1$ approximates
the $W_2$-geodesic (displacement interpolation) between the noise
and data distributions.  OT-CFM~\citep{tong2023ot} aligns noise to
data via a mini-batch optimal transport coupling, further straightening
the learned flow and reducing the number of integration steps needed.
\end{advancedmath}

\textbf{Training recipe.}  Sample a time $t\sim\mathrm{Uniform}(0,1)$,
a noise vector $x_0\sim\calN(0,I)$, and a data sample $x_1$ from the
dataset.  Form the interpolant $x_t=(1-t)x_0+tx_1$.  Train the
network to predict the velocity $x_1-x_0$ from $x_t$ and $t$.

\textbf{Sampling recipe.}  Start from $x_0\sim\calN(0,I)$ and solve
$\dot x_t=v_\theta(x_t,t)$ forward from $t=0$ to $t=1$ using any ODE
solver (e.g., Euler, RK45).  The endpoint $x_1$ is the generated
sample.

\begin{algorithm}[ht]
\caption{Conditional Flow Matching (CFM) --- Training}\label{alg:cfm-train}
\begin{algorithmic}[1]
\Require Dataset $\mathcal{D}$, velocity network $v_\theta$
\Repeat
  \State Sample $x_1\sim\mathcal{D}$ (data), $x_0\sim\calN(0,I)$ (noise), $t\sim\mathrm{Uniform}(0,1)$
  \State $x_t\gets(1-t)\,x_0+t\,x_1$ \Comment{Linear interpolant}
  \State $v^*\gets x_1-x_0$ \Comment{Target velocity}
  \State Take gradient step on $\norm{v_\theta(x_t,t)-v^*}^2$
\Until{converged}
\end{algorithmic}
\end{algorithm}

\begin{algorithm}[ht]
\caption{Flow Matching --- Sampling (Euler method)}\label{alg:fm-sample}
\begin{algorithmic}[1]
\Require Trained $v_\theta$, number of steps $N$, step size $\Delta t=1/N$
\State $x_0\sim\calN(0,I)$
\For{$n=0,1,\ldots,N-1$}
  \State $t\gets n\cdot\Delta t$
  \State $x_{t+\Delta t}\gets x_t+\Delta t\cdot v_\theta(x_t,t)$ \Comment{Euler step along velocity field}
\EndFor
\State \Return $x_1$ \Comment{Generated sample}
\end{algorithmic}
\end{algorithm}

\begin{algorithm}[ht]
\caption{Guided Flow Matching --- Sampling with Score Correction}\label{alg:guided-fm}
\begin{algorithmic}[1]
\Require Trained $v_\theta$, guidance score $s_\mathrm{guide}(\cdot,t)$, weight $\lambda$, steps $N$
\State $x_0\sim\calN(0,I)$
\For{$n=0,1,\ldots,N-1$}
  \State $t\gets n/N$
  \State $v_\mathrm{guided}\gets v_\theta(x_t,t)+\lambda\cdot s_\mathrm{guide}(x_t,t)$ \Comment{Fused velocity}
  \State $x_{t+1/N}\gets x_t+\frac{1}{N}\cdot v_\mathrm{guided}$
\EndFor
\State \Return $x_1$
\end{algorithmic}
\end{algorithm}

\subsection{Why flow matching is preferred for robotics}

For robot trajectory planning, flow matching offers several specific
advantages over diffusion:
\begin{enumerate}[nosep]
  \item \textbf{Speed:} real-time planning requires generating
    trajectories in milliseconds; the 5--10$\times$ speedup of ODE-based
    sampling over iterative denoising is critical for deployment.
  \item \textbf{Determinism:} the ODE sampler is deterministic given an
    initial noise vector, making it easier to analyze, debug, and
    reproduce trajectories.
  \item \textbf{Smooth trajectories:} the straight-line interpolant
    naturally produces smooth paths in trajectory space, which is
    desirable for robot motion.
  \item \textbf{Easy guidance:} modifying the velocity field by adding a
    correction term $v_\theta+u_t$ is straightforward and connects
    directly to optimal control theory (\Cref{sec:pmp}).
\end{enumerate}

\begin{advancedmath}
\subsection{Flow matching and scores are interconvertible}

If the same $p_t$ is generated by velocity $v_t$ and SDE drift $b_t$
with diffusion $g_t$:
$\grad\log p_t=\frac{2}{g_t^2}[v_t-b_t]$.
This means flow matching and score-based diffusion are two views of the
same underlying density evolution---any result about scores can be
translated to velocities and vice versa.
\end{advancedmath}
%==============================================================================
\section{Stochastic Optimal Control}\label{sec:soc}
%==============================================================================

Stochastic optimal control (SOC) connects generative modeling to
control theory. The key insight: computing the optimal control for a
diffusion process is equivalent to computing a score function via
Monte Carlo sampling.  This is why MBD's importance-weighted rollout
strategy works.

\begin{advancedmath}

\subsection{Setup}

Controlled SDE: $\dd X_t=[b(X_t)+\sigma(X_t)u_t]\dd t+\sigma(X_t)\dd W_t$,
cost $J=\E[\int(\Psi+\frac{1}{2}\norm{u}^2)\dd t+h(X_T)]$.

HJB: $-\partial_tV+\Psi+b^\top\grad V+\frac{1}{2}\tr(\sigma\sigma^\top\nabla^2V)
-\frac{1}{2}\norm{\sigma^\top\grad V}^2=0$,
optimal control: $u^*=-\sigma^\top\grad V$.

\subsection{Kappen's log-transform}

Set $\phi=\exp(-V)$. Non-linear HJB becomes \emph{linear}:
$-\partial_t\phi+\Psi\phi+b^\top\grad\phi+\frac{1}{2}\tr(\sigma\sigma^\top\nabla^2\phi)=0$.

\begin{proof}[Derivation]
Substitute $V=-\log\phi$ into the HJB. Using
$\grad V=-\grad\phi/\phi$ and
$\tfrac12\tr(\sigma\sigma^\top\nabla^2 V)
=-\tfrac12\tr(\sigma\sigma^\top\nabla^2\phi)/\phi+\tfrac12\norm{\sigma^\top\grad\phi}^2/\phi^2$,
the quadratic gradient term $-\tfrac12\norm{\sigma^\top\grad V}^2$
cancels, producing the linear PDE.
\end{proof}

By Feynman--Kac:
$\phi(t,x)=\E^{\mathbb{Q}_{t,x}}[\exp(-h(X_T)-\int\Psi\dd s)]$,
so $V=-\log\phi$ is a path integral over uncontrolled rollouts.

\textbf{KL-control / Donsker--Varadhan identity.}
Let $\mathbb{Q}$ be the uncontrolled path measure and $\mathbb{P}^u$
the controlled one. Then
\[
  V(0,x)=\inf_{\mathbb{P}^u\ll\mathbb{Q}}\Bigl\{\E^{\mathbb{P}^u}
  \bigl[h(X_T)+\!\!\int_0^T\!\Psi\,\dd s\bigr]+\KL(\mathbb{P}^u\|\mathbb{Q})\Bigr\},
\]
with optimizer $\dd\mathbb{P}^{u^*}/\dd\mathbb{Q}\propto
\exp(-h(X_T)-\int_0^T\Psi\,\dd s)$.
This is the variational identity underlying MBD and SOCM: the target
trajectory distribution is the optimally-tilted path measure, and its
score is the gradient of the log of this Feynman--Kac expectation.

This is the justification for MBD's Monte Carlo score estimator.

\subsection{SOCM: the score--control bridge}

\citet{domingoenrich2024socm} show that the optimal control has the
form $u^*(t,x)=\E^{\mathbb{Q}_{t,x}}[Z\,\xi_{t,T}]$ where
$Z=\exp(-h(X_T)-\int_t^T\Psi\,\dd s)/\phi(t,x)$ is an explicit
path-functional weight and $\xi_{t,T}$ is a stochastic integral of
Brownian increments. Consequently $u^*$ minimizes
$\min_\theta\E_{t,x}\norm{u_\theta(t,x)-Z\,\xi_{t,T}}^2$
over \emph{uncontrolled} rollouts---structurally identical to DSM
(\Cref{thm:dsm}).
The optimal control is a conditional expectation of weighted path
functionals, just as the optimal score is a conditional expectation of
$\grad\log p_{t|0}$. This is the cleanest formal bridge between score
matching and stochastic optimal control, and it justifies estimating
$\smodel$ by Monte Carlo rollouts of the system dynamics.
\end{advancedmath}

%==============================================================================
\section{Product-of-Experts Guidance}\label{sec:poe}
%==============================================================================

\subsection{Classifier guidance}

From Bayes' rule:
$\grad\log p(x|y)=\grad\log p(x)+\grad\log p(y|x)$.
Posterior score = prior score + likelihood score. This is additive in
scores, multiplicative in densities \citep{dhariwal2021diffusion}.

\begin{advancedmath}
\subsection{Guidance as regularized optimization}

The minimizer of $\E_q[J]+\beta^{-1}\KL(q\|p)$ is
$q^*\propto p\exp(-\beta J)$ with score
$\grad\log q^*=\grad\log p-\beta\grad J$.
Temperature $\beta$ trades off prior vs.\ cost.  This places
classifier guidance in the same variational framework as the
KL-control identity (\Cref{sec:soc}): guidance is entropy-regularized
policy optimization on the data distribution.

\subsection{CFG and policy improvement}

CFG: $\tilde s=(1+w)s_\theta(x,t,y)-ws_\theta(x,t,\varnothing)$.
\citet{cfgrl} show this is a policy improvement operator: the induced
policy $\pi(a|s)\propto\hat\pi(a|s)\exp(wA(s,a))$ has expected return
monotone in $w$.
\end{advancedmath}

\subsection{Product-of-experts for trajectories}

Our fusion \eqref{eq:poe-intro} is the score of
$p^{\mathrm{target}}\propto p^{\mathrm{data}}\cdot[p^{\mathrm{mod}}]^\lambda$
where $p^{\mathrm{mod}}\propto\exp(-J/\lambda)$.
$\lambda\to\infty$: pure optimization. $\lambda\to 0$: behavioral cloning.

\begin{algorithm}[ht]
\caption{Product-of-Experts Fused Trajectory Generation (Ours)}\label{alg:poe-fused}
\begin{algorithmic}[1]
\Require Trained data score $s_\theta$, dynamics $f$, cost $J$, temperature $\lambda$,
  adaptive weight $\lambda(x,t)$, MC samples $M$, steps $N$
\State $\tau_0\sim\calN(0,I)$ \Comment{Start from noise (flow-matching convention)}
\For{$n=0,1,\ldots,N-1$}
  \State $t\gets n/N$
  \State $s_\mathrm{data}\gets s_\theta(\tau_t,t)$
    \Comment{Data-driven score (from trained flow/diffusion model)}
  \State $s_\mathrm{mod}\gets\textsc{MBD-Score}(\tau_t,f,J,\lambda,t,M)$
    \Comment{Model-based score (\Cref{alg:mbd-score})}
  \State $s_\mathrm{fused}\gets s_\mathrm{data}+\lambda(\tau_t,t)\cdot s_\mathrm{mod}$
    \Comment{Adaptive fusion}
  \State $\tau_{t+1/N}\gets\tau_t+\frac{1}{N}\cdot v(s_\mathrm{fused},\tau_t,t)$
    \Comment{ODE step with fused velocity}
\EndFor
\State \Return $\tau_1$ \Comment{Fused trajectory: diverse, feasible, low-cost}
\end{algorithmic}
\end{algorithm}

%==============================================================================
\section{Model-Based Diffusion}\label{sec:mbd}
%==============================================================================

Model-Based Diffusion (MBD)~\citep{mbd} is the direct precursor of
the model-based score in our project.  The key insight is that for
systems with known dynamics, we do \emph{not} need a neural network to
compute the score---we can compute it analytically from the dynamics
model and cost function using Monte Carlo sampling.

\textbf{How MBD works (overview).}
\begin{enumerate}[nosep]
  \item Define a Boltzmann trajectory distribution
    $p^{\mathrm{mod}}(\tau)\propto\exp(-J(\tau)/\lambda)$ that assigns
    high probability to low-cost, dynamically feasible trajectories.
  \item Noise this distribution via a variance-preserving SDE (the same
    noising used in DDPM) to get noised densities $p^{\mathrm{mod}}_t$.
  \item At each denoising step, compute the score
    $\grad\log p^{\mathrm{mod}}_t$ via Tweedie's formula and
    importance-weighted sampling over dynamically feasible rollouts.
  \item Use this score in the standard reverse SDE to generate
    trajectories.
\end{enumerate}

MBD solves Maze2D, Block Push, and Franka manipulation benchmarks
\emph{without any demonstration data}, matching or outperforming
Diffuser~\citep{janner2022} (which requires offline data).

\begin{advancedmath}
\subsection{Tweedie's formula}

If $x_t=\alpha_tx_0+\sigma_t\epsilon$, $\epsilon\sim\calN(0,I)$:
\begin{equation}
  \grad\log p_t(x_t)=\frac{\alpha_t}{\sigma_t^2}(\E[x_0|x_t]-x_t/\alpha_t).
  \label{eq:tweedie}
\end{equation}
The score points from noisy $x_t$ toward the expected clean $x_0$.
This is the computational workhorse of MBD: if we can estimate the
posterior mean $\E[x_0|x_t]$, we immediately have the score.

\begin{proof}
$\grad\log p_t(x_t)=\E[\grad_{x_t}\log p_{t|0}(x_t|X_0)\mid X_t=x_t]$.
Substituting $\grad\log p_{t|0}=-\sigma_t^{-2}(x_t-\alpha_tx_0)$ gives
$\grad\log p_t=(\alpha_t/\sigma_t^2)(\E[X_0|X_t]-x_t/\alpha_t)$.
\end{proof}

\subsection{MBD's Monte Carlo estimator}

Estimate $\E[x_0|x_t]$ via importance sampling:
$\widehat\E[x_0|x_t]=\sum w_i\tau_0^{(i)}/\sum w_i$,
$w_i=\exp(-J(\tau_0^{(i)})/\lambda)$. This is ``importance-weighted MPC
in score space.''

\begin{proposition}[CLT for self-normalized IS]
Assume $\E_q[w^2\tau_0^2]<\infty$ and $\E_q[w]>0$. Then
$\sqrt{M}(\widehat\E[X_0|X_t]-\E[X_0|X_t])\to\calN(0,\Sigma)$ in
distribution, with $\Sigma=\Var_q[w(\tau_0-\E[X_0|X_t])]/\E_q[w]^2$.
In particular the MBD score estimator via
\eqref{eq:tweedie} is consistent at rate $M^{-1/2}$.
\end{proposition}

As $\lambda\to\infty$, the Boltzmann distribution concentrates on
$\argmin J$ and the reverse SDE reduces to smoothed, annealed
gradient descent on $J$---exactly sampling-based MPC.
\end{advancedmath}

\begin{algorithm}[ht]
\caption{MBD: Model-Based Score Estimation at Denoising Step $k$}\label{alg:mbd-score}
\begin{algorithmic}[1]
\Require Noised trajectory $\tau_k$, dynamics $f$, cost $J$, temperature $\lambda$,
  noise schedule $\bar\alpha_k$, number of MC samples $M$
\State $\mu_\mathrm{post}\gets\tau_k/\sqrt{\bar\alpha_k}$,\quad
  $\sigma^2_\mathrm{post}\gets(1-\bar\alpha_k)/\bar\alpha_k$
  \Comment{Gaussian posterior $q(\tau_0|\tau_k)$}
\For{$i=1,\ldots,M$}
  \State Sample $\tau_0^{(i)}\sim\calN(\mu_\mathrm{post},\sigma^2_\mathrm{post}I)$
  \State Project $\tau_0^{(i)}$ onto dynamically feasible trajectories via $f$
    \Comment{Forward rollout}
  \State $w_i\gets\exp(-J(\tau_0^{(i)})/\lambda)$ \Comment{Boltzmann weight}
\EndFor
\State $\widehat\E[\tau_0|\tau_k]\gets\sum_{i=1}^M w_i\tau_0^{(i)}\big/\sum_{i=1}^M w_i$
  \Comment{Self-normalized IS}
\State $\hat s_\mathrm{mod}\gets\frac{\sqrt{\bar\alpha_k}}{1-\bar\alpha_k}
  \bigl(\widehat\E[\tau_0|\tau_k]-\tau_k/\sqrt{\bar\alpha_k}\bigr)$
  \Comment{Tweedie's formula}
\State \Return $\hat s_\mathrm{mod}$
\end{algorithmic}
\end{algorithm}

\begin{algorithm}[ht]
\caption{MBD: Full Trajectory Generation (No Training Data)}\label{alg:mbd-full}
\begin{algorithmic}[1]
\Require Dynamics $f$, cost $J$, temperature $\lambda$, MC samples $M$, denoising steps $T$
\State $\tau_T\sim\calN(0,I)$ \Comment{Start from noise}
\For{$k=T,T-1,\ldots,1$}
  \State $\hat s_\mathrm{mod}\gets\textsc{MBD-Score}(\tau_k,f,J,\lambda,\bar\alpha_k,M)$
    \Comment{\Cref{alg:mbd-score}}
  \State $\mu_k\gets\frac{1}{\sqrt{\alpha_k}}\bigl(\tau_k+(1-\alpha_k)\hat s_\mathrm{mod}\bigr)$
    \Comment{Reverse SDE mean}
  \State $\tau_{k-1}\sim\calN(\mu_k,\sigma_k^2 I)$
\EndFor
\State \Return $\tau_0$ \Comment{Feasible, low-cost trajectory}
\end{algorithmic}
\end{algorithm}

\subsection{Limitations of MBD}

MBD relies entirely on the analytical model and therefore has several
important limitations that motivate our project:
\begin{itemize}[nosep]
  \item \textbf{No multimodality:} MBD cannot capture
    demonstration-informed behaviors---it produces feasible but
    potentially unnatural trajectories (e.g., ignoring preferred grasps
    or learned obstacle-avoidance heuristics).
  \item \textbf{Model dependence:} it assumes a sufficiently accurate
    dynamics model, which may fail under model mismatch, unmodeled
    contacts, or sim-to-real transfer.
  \item \textbf{Variance:} the importance-sampling score estimator
    suffers high variance when $\lambda$ is small (peaked Boltzmann) or
    the trajectory dimension is large.
  \item \textbf{No safety guarantees:} MBD provides no formal
    collision-avoidance guarantees.
\end{itemize}
Our work targets these limitations: the learned prior restores
distributional richness, the adaptive $\lambda(x,t)$ weighting remains
robust under model imperfection, and the data score anchors the
denoising process to reduce variance.  We additionally plan to
incorporate a CBF-derived score term~\citep{safeflow} for safety.

%==============================================================================
\section{The LQ--Riccati--Score Bridge}\label{sec:lqbridge}
%==============================================================================

\begin{theorem}[LQ model-score = Riccati feedback]\label{thm:lq-score}
Under LQR with $\Sigma_w=0$, define
$p^{\mathrm{mod}}(\tau|x_0)\propto\exp(-J(\tau)/\lambda)$. Then at the
optimal trajectory:
\begin{equation}
  \boxed{\grad_{x_k}\log p^{\mathrm{mod}}(\tau|x_0)\big|_{\tau^*}
  =-\frac{1}{\lambda}P_kx_k^*}
  \label{eq:lq-score}
\end{equation}
\end{theorem}

\textbf{Interpretation.}
The fused drift becomes $s_\theta(\tau,t)-P_kx_k$: learned data score
+ classical LQR feedback. As $\lambda\to\infty$, we get pure LQR.
As $\lambda\to 0$, pure behavioral cloning. Modern generative models
\emph{contain} classical optimal control as a limiting case.
If the data distribution is also Gaussian, the fused process is
linear-Gaussian and exactly solvable (our analytic testbed).

\begin{advancedmath}
\begin{proof}[Full derivation of \Cref{thm:lq-score}]
\textbf{Step 1 (Quadratic cost in control space).}
Under $x_{k+1}=Ax_k+Bu_k$, unrolling gives $x_k=\Phi_kx_0+\Gamma_ku$
with $\Phi_k,\Gamma_k$ linear operators. Substituting into
$J=\sum x_k^\top Qx_k+u_k^\top Ru_k+x_H^\top Q_Hx_H$ yields
$J(\tau)=u^\top Mu+2u^\top Nx_0+x_0^\top Lx_0$ with
$M\succ 0$. Therefore
$p^{\mathrm{mod}}(u|x_0)=\calN(-M^{-1}Nx_0,\;\lambda M^{-1})$:
a Gaussian centered at the LQR solution.

\textbf{Step 2 (Bellman factorization).}
Fix $k$ and condition on $x_{<k}$. The cost-to-go
$J_{k\to H}(x_k,u_{k:H-1})$ decomposes as
$x_k^\top P_kx_k+\delta_k^\top M_k\delta_k$
where $\delta_k=u_{k:H-1}-u^*_{k:H-1}(x_k)$ and $M_k\succ 0$
(Riccati Hessian). No linear cross term exists because
$u^*_{k:H-1}$ is the unconstrained minimizer. Hence
$\int\exp(-J_{k\to H}/\lambda)\,\dd u_{k:H-1}
\propto\exp(-x_k^\top P_kx_k/\lambda)$.

\textbf{Step 3 (Score computation).}
$p^{\mathrm{mod}}(x_k|x_{<k})\propto\exp(-x_k^\top P_kx_k/\lambda)$
as a function of $x_k$, so
$\grad_{x_k}\log p^{\mathrm{mod}}=-\lambda^{-1}\cdot 2P_kx_k$.
Under the convention $J=\tfrac12\sum(x^\top Qx+u^\top Ru)$ the
factor of 2 is absorbed, giving \eqref{eq:lq-score}.
\end{proof}

\begin{corollary}[Linear-Gaussian fused ODE]
Under a Gaussian data distribution
$p_t^{\mathrm{data}}(\tau)=\calN(\mu_t,\Sigma_t)$ (straight-line CFM)
and the LQ model distribution, the product-of-experts drift is affine:
$s_\theta+\lambda\smodel=-\Sigma_t^{-1}(\tau-\mu_t)-P_\tau\tau$,
where $P_\tau=\mathrm{blkdiag}(P_0,\ldots,P_H)$.
The fused generative process is itself linear-Gaussian and exactly
solvable: its trajectory is the LQG-smoothed MAP estimate.
\end{corollary}
\end{advancedmath}

%==============================================================================
\section{Generative Models for Robot Trajectory Planning}\label{sec:robotplan}
%==============================================================================

This section introduces the key algorithms that apply diffusion and
flow matching to robot trajectory planning.  Each method generates
trajectories by sampling from a learned distribution, but they differ
in architecture, conditioning mechanism, and what knowledge (data,
dynamics, or both) they exploit.

\subsection{Diffuser: the first diffusion-based planner}

Diffuser~\citep{janner2022} was the first to use a denoising diffusion
model for long-horizon planning.  It trains a DDPM over entire
trajectories $\tau=(s_0,a_0,s_1,a_1,\ldots,s_H)$ from an offline
dataset.  At test time, the start state is enforced by
\emph{inpainting} (clamping known values at each denoising step), and
reward maximization uses classifier guidance: the denoising mean is
shifted by the gradient of a learned return estimator.  Diffuser
matches or outperforms Decision Transformer and CQL on D4RL benchmarks
while naturally handling goal-conditioned and reward-maximizing
settings with the same model.

\textbf{Relevance:} Our project generalizes Diffuser's guidance
mechanism from a learned reward gradient to an analytically computed
model-based score derived from dynamics and cost.

\begin{algorithm}[ht]
\caption{Diffuser: Classifier-Guided Trajectory Planning}\label{alg:diffuser}
\begin{algorithmic}[1]
\Require Trained trajectory denoiser $\epsilon_\theta$, learned return estimator $\mathcal{J}$,
  guidance weight $w$, start state $s_0$, noise schedule $\{\bar\alpha_k,\sigma_k\}$
\State $\tau_T\sim\calN(0,I)$ \Comment{Sample noise trajectory}
\For{$k=T,T-1,\ldots,1$}
  \State $\mu_k\gets\frac{1}{\sqrt{\alpha_k}}\bigl(\tau_k-\frac{1-\alpha_k}{\sqrt{1-\bar\alpha_k}}\epsilon_\theta(\tau_k,k)\bigr)$
    \Comment{Predicted mean}
  \State $\tilde\mu_k\gets\mu_k+w\cdot\sigma_k^2\cdot\grad_{\tau_k}\mathcal{J}(\mu_k)$
    \Comment{Classifier guidance}
  \State $\tau_{k-1}\sim\calN(\tilde\mu_k,\sigma_k^2 I)$
  \State $\tau_{k-1}[s_0]\gets s_0$ \Comment{Inpainting: clamp start state}
\EndFor
\State \Return $\tau_0$ \Comment{Planned trajectory}
\end{algorithmic}
\end{algorithm}

\subsection{Diffusion Policy: action-chunk diffusion for manipulation}

Diffusion Policy~\citep{chi2023diffusionpolicy} applies DDPM to
generate \emph{action chunks} $(a_t,\ldots,a_{t+H_a})$ conditioned on
the current observation $o_t$ (visual or low-dimensional).  It uses a
1D temporal U-Net architecture, treating the action horizon as the
spatial dimension, with observations injected via FiLM conditioning or
cross-attention.  A receding-horizon scheme applies the first
$H_e<H_a$ actions and re-plans.

Diffusion Policy achieves state-of-the-art on 11/12 RoboMimic and
Push-T tasks, including real-robot bimanual manipulation, and captures
multimodal behavior (e.g., going left or right around an obstacle) that
deterministic policies and VAEs miss.

\textbf{Relevance:} Our data-driven score $s_\theta$ is trained in the
same framework; we add the model-based correction on top.

\begin{algorithm}[ht]
\caption{Diffusion Policy: Receding-Horizon Action Generation}\label{alg:diffpolicy}
\begin{algorithmic}[1]
\Require Trained conditional denoiser $\epsilon_\theta(\cdot,\cdot|o)$, observation $o_t$,
  action horizon $H_a$, execution horizon $H_e<H_a$
\State $a^T_{t:t+H_a}\sim\calN(0,I)$ \Comment{Initialize noise action chunk}
\For{$k=T,T-1,\ldots,1$}
  \State $a^{k-1}_{t:t+H_a}\gets\mathrm{DDPMStep}(\epsilon_\theta,a^k_{t:t+H_a},k,o_t)$
    \Comment{Conditional denoising step}
\EndFor
\State Execute $a^0_{t},a^0_{t+1},\ldots,a^0_{t+H_e-1}$ on robot \Comment{Execute first $H_e$ actions}
\State Observe new $o_{t+H_e}$; repeat from line 1 \Comment{Receding horizon}
\end{algorithmic}
\end{algorithm}

\subsection{MPD: cost-guided motion planning diffusion}

Motion Planning Diffusion (MPD)~\citep{carvalho2023mpd} trains a
trajectory diffusion model on collision-free paths, then at test time
composes the score with cost-function gradients:
$\tilde s=s_\theta+\lambda\grad c(\hat\tau_0)$,
where $\hat\tau_0$ is the Tweedie-denoised ``clean'' trajectory
estimate and $c$ includes collision, smoothness, and goal costs.

\textbf{Relevance:} MPD's cost-guidance mechanism is a predecessor of
our model-based score.  The difference: MPD uses ad hoc cost gradients
on the Tweedie estimate; we compute $\smodel$ from the full Boltzmann
distribution and analyze it rigorously.

\subsection{FlowMP: flow matching for robot motion}

FlowMP~\citep{flowmp} is one of the first methods to apply conditional
flow matching to robot motion planning.  It trains a CFM model on
demonstration trajectories represented as sequences of joint
configurations, conditioned on context (start, goal, obstacles).  A key
feature is the second-order formulation: the velocity field operates on
a phase-space representation $(\tau,\dot\tau)$, encouraging kinematic
feasibility.

FlowMP achieves 5--10$\times$ faster inference than Diffusion Policy
and MPD with comparable success rates on Franka manipulation benchmarks,
validating flow matching as a practical alternative to DDPM for robot
planning.

\textbf{Relevance:} Our project uses the same CFM backbone and augments
it with model-based score guidance.

\subsection{GoalFlow: one-step flow matching for driving}

GoalFlow~\citep{xing2025goalflow} applies conditional flow matching to
end-to-end autonomous driving with goal-conditioned trajectory
generation.  Its key innovation is one-step distillation: a student
network is trained to produce high-quality trajectories in a single
forward pass (rather than 10--20 ODE steps), enabling real-time
deployment at under 50ms per trajectory.

\textbf{Relevance:} GoalFlow shows flow matching scales to
high-dimensional driving scenarios.  The one-step distillation idea
could accelerate our fused generative process at deployment time.

\subsection{BSD: training-free kernel-based scores}

Behavioral Score Diffusion (BSD)~\citep{bsd} computes the trajectory
score directly from a dataset using kernel density estimation (KDE),
requiring no neural network training.  Given a dataset
$\{\tau^{(i)}\}_{i=1}^N$, the score at noise level $t$ is available in
closed form as a weighted sum over data points---essentially a soft
nearest-neighbor regression in the noised space.

\textbf{Relevance:} BSD provides a training-free ``behavioral score''
$s_\text{data}$ that can be directly fused with MBD's model-based score
via the product-of-experts mechanism.  This is the simplest realization
of our score-fusion framework.

\subsection{OC-Flow: optimal-control guidance for flow matching}

OC-Flow~\citep{ocflow} guides a pre-trained flow-matching model at
test time to maximize a reward $R(x_1)$ without retraining.  It views
the flow ODE $\dot x=v_\theta+u_t$ as a deterministic optimal control
problem and solves it via the Pontryagin Maximum Principle: a
forward-backward sweep computes the adjoint $\lambda_t$ and sets
$u_t=-\lambda_t$.

\textbf{Relevance:} OC-Flow shows how PMP provides a principled
guidance mechanism for flow matching.  Our project uses a different
(score-based) guidance mechanism, but OC-Flow's adjoint framework could
serve as an alternative or complement.

\begin{algorithm}[ht]
\caption{OC-Flow: PMP-Based Guidance for Flow Matching}\label{alg:ocflow}
\begin{algorithmic}[1]
\Require Trained $v_\theta$, reward $R$, steps $N$, number of iterations $K$
\State $x_0\sim\calN(0,I)$; $u_t^{(0)}\gets 0\;\forall t$ \Comment{Initialize}
\For{$\mathrm{iter}=1,\ldots,K$}
  \State \textbf{Forward pass:} integrate $\dot x_t=v_\theta(x_t,t)+u_t^{(\mathrm{iter}-1)}$
    from $t=0$ to $t=1$
  \State Compute $\grad R(x_1)$ \Comment{Reward gradient at endpoint}
  \State \textbf{Backward pass:} set $\lambda_1=-\grad R(x_1)$; integrate
    $\dot\lambda_t=-(\partial_x v_\theta)^\top\lambda_t$ backward from $t=1$ to $t=0$
  \State $u_t^{(\mathrm{iter})}\gets-\lambda_t\;\forall t$ \Comment{Optimal control from costate}
\EndFor
\State Integrate $\dot x_t=v_\theta(x_t,t)+u_t^{(K)}$ from $t=0$ to $t=1$
\State \Return $x_1$ \Comment{Reward-guided sample}
\end{algorithmic}
\end{algorithm}

\subsection{SOCM: the score--control bridge}

Stochastic Optimal Control Matching (SOCM)~\citep{domingoenrich2024socm}
solves high-dimensional stochastic optimal control problems by
expressing the optimal control as a conditional expectation under the
uncontrolled diffusion---structurally identical to denoising score
matching.  This provides the formal bridge between score matching and
optimal control: computing model-based scores via Monte Carlo rollouts
(exactly MBD's strategy) is justified by the Feynman--Kac identity.

\subsection{CFGRL: guidance as policy improvement}

\citet{cfgrl} provide a theoretical analysis of classifier-free
guidance (CFG) as a policy improvement operator.  CFG with weight $w$
effectively samples from a tilted policy
$\pi_w(a|s)\propto\pi_\theta(a|s)\cdot(\pi_\theta(a|s)/\pi_\theta(a))^w$,
which upweights actions most relevant to the current state.  They prove
a soft policy improvement guarantee for sufficiently small $w$, with
mode collapse for large $w$---predicting the
$\lambda$-sensitivity we hypothesize in our framework.

\subsection{DIAL-MPC: diffusion-style annealing for MPC}

DIAL-MPC~\citep{dialmpc} applies diffusion-style annealed sampling to
model predictive control for torque-level locomotion.  Rather than using
a learned score network, it uses the MPC cost function directly within
a multi-step annealing scheme inspired by the denoising process.  This
demonstrates that the diffusion framework can enhance classical
sampling-based MPC even without a learned component.

%==============================================================================
\section{Safety in Generative Planning}\label{sec:safety}
%==============================================================================

A major limitation of unconstrained generative planners is that they
provide no formal safety guarantees.  Several recent methods address
this by incorporating \emph{control barrier functions} (CBFs) or
\emph{Lyapunov functions} into the generative process.

\textbf{Control barrier functions (CBFs).}
A CBF defines a safe set $\mathcal{C}=\{x:h(x)\ge 0\}$ and enforces
forward invariance via the condition
$\dot h(x,u)\ge-\alpha(h(x))$ for a class-$\mathcal{K}$ function
$\alpha$.  Applied to a velocity field $v$, this becomes a linear
constraint on $v$ (since $\dot h=\grad h^\top v$), yielding a
quadratic program (QP) at each integration step.

\subsection{SafeFlow}

SafeFlow~\citep{safeflow} projects the flow-matching velocity through a
CBF safety filter at each ODE step:
$v_\text{safe}=\argmin_v\norm{v-v_\theta}^2$ subject to the CBF
constraint.  This achieves zero constraint violations on navigation and
Franka planning benchmarks.

\begin{algorithm}[ht]
\caption{SafeFlow: CBF-Filtered Flow Matching Sampling}\label{alg:safeflow}
\begin{algorithmic}[1]
\Require Trained $v_\theta$, barrier function $h$ (safe set: $h(x)\ge 0$),
  class-$\mathcal{K}$ function $\alpha$, steps $N$
\State $x_0\sim\calN(0,I)$
\For{$n=0,1,\ldots,N-1$}
  \State $t\gets n/N$
  \State $v_\mathrm{nom}\gets v_\theta(x_t,t)$ \Comment{Nominal velocity from flow model}
  \If{$\grad h(x_t)^\top v_\mathrm{nom}<-\alpha(h(x_t))$} \Comment{CBF violated?}
    \State Solve QP: $v_\mathrm{safe}\gets\argmin_{v}\norm{v-v_\mathrm{nom}}^2$
      \textbf{ s.t.} $\grad h(x_t)^\top v\ge-\alpha(h(x_t))$
  \Else
    \State $v_\mathrm{safe}\gets v_\mathrm{nom}$ \Comment{Already safe}
  \EndIf
  \State $x_{t+1/N}\gets x_t+\frac{1}{N}\cdot v_\mathrm{safe}$
\EndFor
\State \Return $x_1$ \Comment{Safe trajectory}
\end{algorithmic}
\end{algorithm}

\subsection{SafeFlowMatcher}

SafeFlowMatcher~\citep{yang2025safeflowmatcher} improves upon
SafeFlow's single-step projection with a predictor-corrector scheme: a
tentative Euler step followed by a CBF-QP correction.  This handles
cases where a single projection is insufficient due to large step sizes
or high curvature of the safe boundary.

\subsection{UniConFlow}

UniConFlow~\citep{yang2025uniconflow} handles multiple heterogeneous
constraints (obstacles, joint limits, kinematic feasibility) in a
unified QP at each ODE step, achieving certified constraint
satisfaction across all constraint types simultaneously.

\subsection{SafeDiffuser}

SafeDiffuser~\citep{xiao2025safediffuser} embeds CBF-based safety
\emph{within} the DDPM denoising loop rather than as a post-hoc filter.
At each denoising step, after the score-based update, a differentiable
CBF correction is applied to the Tweedie estimate.  A control Lyapunov
function (CLF) additionally ensures progress toward the goal, yielding
a CBF-CLF-QP at each step.

\subsection{$S^2$Diff: Lyapunov-guided diffusion}

$S^2$Diff~\citep{s2diff} constructs a score function from a known
Lyapunov function $V(x)$ for the system $\dot x=f(x)+g(x)u$, pointing
in the direction that decreases the Lyapunov function most rapidly.
This provides provably stable control on nonlinear systems without any
demonstration data.

\textbf{Relevance:} $S^2$Diff is the closest relative to our
model-based score concept: both derive an analytical score from system
properties.  $S^2$Diff uses a Lyapunov function (stability) while we
use a cost function (optimality).  In the LQ case, both reduce to the
Riccati matrix.

\subsection{Safety in our framework}

In our product-of-experts framework, a CBF constraint can be absorbed
as an additional score term:
$s_\text{CBF}=\grad\log\mathbf{1}[h(\tau)\ge 0]$, or more practically,
as a post-processing QP layer (SafeFlow-style) applied after the fused
score update.  We plan to incorporate SafeFlow's CBF mechanism as an
independent safety module.

%==============================================================================
\section{Robotics Survey}\label{sec:survey}
%==============================================================================

\begin{table}[ht]
  \centering\small
  \caption{Trajectory methods by score source.}
  \label{tab:survey}
  \begin{tabular}{p{3cm}cccp{3.8cm}}
    \toprule
    \textbf{Method} & \textbf{Year} & \textbf{Data} & \textbf{Model} & \textbf{Formalism}\\
    \midrule
    Diffuser \citep{janner2022} & 2022 & \checkmark & $\times$ & classifier guidance\\
    Diff.\ Policy \citep{chi2023diffusionpolicy} & 2023 & \checkmark & $\times$ & DDPM actions\\
    MPD \citep{carvalho2023mpd} & 2023 & \checkmark & partial & cost guidance\\
    OC-Flow \citep{ocflow} & 2025 & \checkmark & reward & PMP adjoint\\
    FlowMP \citep{flowmp} & 2025 & \checkmark & $\times$ & second-order CFM\\
    SafeFlow \citep{safeflow} & 2025 & \checkmark & safety & CBF-guided FM\\
    SafeFlowMatcher \citep{yang2025safeflowmatcher} & 2025 & \checkmark & safety & CBF-QP\\
    UniConFlow \citep{yang2025uniconflow} & 2025 & \checkmark & both & constrained FM\\
    GoalFlow \citep{xing2025goalflow} & 2025 & \checkmark & $\times$ & one-step CFM\\
    SafeDiffuser \citep{xiao2025safediffuser} & 2025 & \checkmark & safety & CBF denoiser\\
    $S^2$Diff \citep{s2diff} & 2025 & $\times$ & \checkmark & Lyapunov score\\
    SOCM \citep{domingoenrich2024socm} & 2024 & --- & \checkmark & SOC matching\\
    MBD \citep{mbd} & 2024 & opt. & \checkmark & MC score\\
    BSD \citep{bsd} & 2026 & \checkmark & $\times$ & kernel score\\
    CFGRL \citep{cfgrl} & 2025 & \checkmark & --- & CFG as PI\\
    DIAL-MPC \citep{dialmpc} & 2025 & $\times$ & \checkmark & annealed sampling MPC\\
    \textbf{Ours} & 2026 & \checkmark & \checkmark & PoE + LQ-Riccati\\
    \bottomrule
  \end{tabular}
\end{table}

No prior work formalizes score-level product-of-experts fusion, studies
the $\lambda$-tradeoff quantitatively, or connects the model-based
correction to the DARE.  Existing methods either use data alone
(Diffuser, Diffusion Policy, FlowMP, BSD), physics alone (MBD,
$S^2$Diff, DIAL-MPC), or add safety post-hoc (SafeFlow, SafeDiffuser).
Our framework unifies all three: a learned data prior, an analytical
model-based score, and a CBF safety layer, combined via principled
score-level fusion with a theoretically grounded connection to
classical optimal control.

%==============================================================================
\section{Open Questions}\label{sec:open}
%==============================================================================

\textbf{Confident:} PoE decomposition is well-posed;
\Cref{thm:lq-score} is analytically verifiable; MBD's estimator is
consistent.

\textbf{Uncertain:} (i) Convergence of fused ODE in the nonlinear case;
(ii) MC score variance in high dimensions; (iii) calibration between
$p^{\mathrm{data}}$ and $p^{\mathrm{mod}}$; (iv) $\lambda$-sensitivity.

\textbf{Fallback:} Frame around (1) \Cref{thm:lq-score}, (2) empirical
$\lambda$-characterization, (3) failure-mode analysis.

%##############################################################################
% APPENDICES: Rigorous Mathematical Reference
%##############################################################################
\appendix
\newpage
\part*{Appendices: Rigorous Mathematical Reference}
\addcontentsline{toc}{part}{Appendices: Rigorous Mathematical Reference}

\noindent\emph{The following appendices reproduce the full
measure-theoretic treatment from the original rigorous literature
review. They are included as a self-contained reference for future
course projects and are not required for understanding the main text.
All random elements live on a complete filtered probability space
$(\Omega,\calF,(\calF_t),\Prob)$ satisfying the usual conditions.}

%==============================================================================
\section{Stochastic Calculus and SDEs}\label{app:stoch}
%==============================================================================

\subsection{Brownian motion and filtrations}

\begin{definition}[$d$-dimensional Brownian motion]\label{def:bm}
A measurable process $W=(W_t)_{t\ge 0}$, $W_t:\Omega\to\R^d$, adapted to
$(\calF_t)$, is an $(\calF_t)$-\emph{Brownian motion} if (i) $W_0=0$ a.s.;
(ii) for $0\le s<t$, $W_t-W_s$ is independent of $\calF_s$ and distributed
as $\calN(0,(t-s)I_d)$; (iii) the paths $t\mapsto W_t(\omega)$ are
continuous for a.e.\ $\omega$.
\end{definition}

\begin{proposition}[Quadratic variation]\label{prop:qv}
For every partition $\pi=\{0=t_0<\dots<t_n=t\}$ of $[0,t]$, with
$\norm{\pi}=\max_i(t_{i+1}-t_i)$,
\[
  \sum_{i=0}^{n-1}(W_{t_{i+1}}-W_{t_i})(W_{t_{i+1}}-W_{t_i})^\top\ \xrightarrow{\,\Prob\,}\ tI_d
  \quad\text{as }\norm{\pi}\to 0.
\]
In particular, $d$-dimensional Brownian paths have infinite total
variation on any interval but finite quadratic variation
$[W]_t=tI_d$, which is the source of the second-order correction in
It\^o's formula.
\end{proposition}

\subsection{The It\^o integral}\label{app:itoint}

Let $\calH^2_T$ denote the Hilbert space of progressively measurable
$\R^{d\times m}$-valued processes $H$ with
$\norm{H}_{\calH^2}^2:=\E\int_0^T\tr(HH^\top)\,\dd t<\infty$.

\begin{definition}[It\^o integral]\label{def:itoint}
For a simple process
$H_t=\sum_{i=0}^{n-1}\xi_i\ind_{(t_i,t_{i+1}]}(t)$, with $\xi_i\in
\calF_{t_i}$ and $\E\norm{\xi_i}^2<\infty$, set
\[
  \int_0^T H_t\,\dd W_t\;:=\;\sum_{i=0}^{n-1}\xi_i\,(W_{t_{i+1}}-W_{t_i}).
\]
The map $H\mapsto\int_0^{\cdot}H_t\dd W_t$ extends by continuity to all of
$\calH^2_T$, yielding a continuous linear isometry into
$L^2(\Omega;C([0,T];\R^d))$:
\begin{equation}
  \E\Bigl\lVert\int_0^T H_t\,\dd W_t\Bigr\rVert^2
  = \E\int_0^T\tr(H_tH_t^\top)\,\dd t
  \qquad(\text{It\^o isometry}).
  \label{eq:ito-isometry}
\end{equation}
\end{definition}

\begin{theorem}[Martingale property]\label{thm:ito-mart}
For $H\in\calH^2_T$ the process $M_t:=\int_0^t H_s\,\dd W_s$ is an
$(\calF_t)$-martingale with continuous paths, $\E M_t=0$, and
$[M]_t=\int_0^t H_sH_s^\top\,\dd s$.
\end{theorem}
\begin{proof}[Proof sketch]
Check these properties for simple integrands (direct computation) and pass
to the $\calH^2_T$-limit using Doob's $L^2$ inequality; see
\citep[Ch.~3]{karatzas1991bm}.
\end{proof}

\begin{remark}[Beyond $\calH^2$]
The integral extends to the class of locally square-integrable progressively
measurable processes by a standard localization argument; the resulting
integral is a continuous local martingale \citep[Ch.~3.2]{karatzas1991bm}.
\end{remark}

\subsection{It\^o's formula}

\begin{theorem}[It\^o--Doeblin formula]\label{thm:itoformula}
Let $X_t$ be an It\^o process
\[
  X_t = X_0 + \int_0^t\mu_s\,\dd s + \int_0^t\Sigma_s\,\dd W_s,
\]
with $\mu$ progressively measurable, $\mu\in L^1(\dd t\otimes\Prob)$, and
$\Sigma\in\calH^2_T$. For every $\varphi\in C^{1,2}([0,T]\times\R^d)$,
\begin{align}
  \varphi(t,X_t)
  &= \varphi(0,X_0)
  + \int_0^t\!\Bigl[\partial_s\varphi(s,X_s)+\inner{\mu_s}{\grad\varphi(s,X_s)}
                +\tfrac12\tr(\Sigma_s\Sigma_s^\top\nabla^2\varphi(s,X_s))\Bigr]\dd s\nonumber\\
  &\quad + \int_0^t\inner{\grad\varphi(s,X_s)}{\Sigma_s\,\dd W_s}.
  \label{eq:itoformula}
\end{align}
\end{theorem}
\begin{proof}[Proof sketch]
Localize so that $\Sigma$ is simple and $\varphi$ has compact support,
apply a second-order Taylor expansion on a partition, and use
\Cref{prop:qv}: the first-order term gives $\mu$ and $\Sigma\dd W$, and the
second-order Taylor term contributes only the trace against
$\dd[X]_s=\Sigma_s\Sigma_s^\top\dd s$; the cross terms vanish by the
$o(\dd t)$ magnitude of $(\dd W_t)^{\otimes 3}$. Pass to the limit in
$L^2$; see \citep[Thm.~3.3.6]{karatzas1991bm}.
\end{proof}

\begin{example}[Log-density computation]\label{ex:dlogp}
If $X_t$ has density $p_t\in C^{1,2}$ and we set
$\varphi(t,x)=\log p_t(x)$, then \eqref{eq:itoformula} gives
\[
  \dd\log p_t(X_t)
  = \Bigl[\partial_t\log p_t+\inner{\mu_t}{\grad\log p_t}
    +\tfrac12\tr(\Sigma_t\Sigma_t^\top\nabla^2\log p_t)\Bigr]\dd t
    +\inner{\grad\log p_t}{\Sigma_t\dd W_t},
\]
which is the fundamental identity behind score-based models: the drift
term reduces under Fokker--Planck to a compensator, and
the Brownian part is the ``score noise'' that denoising score matching
targets.
\end{example}

\subsection{SDEs: existence and uniqueness}

\begin{definition}[Strong solution]\label{def:strong}
Given Borel-measurable $b:[0,T]\times\R^d\to\R^d$,
$\sigma:[0,T]\times\R^d\to\R^{d\times m}$, an initial law $p_0$, and a
fixed $(\calF_t)$-Brownian motion $W$, a \emph{strong solution} of
\begin{equation}
  \dd X_t = b(t,X_t)\,\dd t + \sigma(t,X_t)\,\dd W_t,\qquad X_0\sim p_0,
  \label{eq:ito-sde}
\end{equation}
is an $(\calF_t)$-adapted continuous process $X$ with
$\int_0^T(\norm{b(t,X_t)}+\norm{\sigma(t,X_t)}^2)\dd t<\infty$ a.s.\ such
that $X_t=X_0+\int_0^t b\,\dd s+\int_0^t\sigma\,\dd W$ for every
$t\in[0,T]$ a.s.
\end{definition}

\begin{theorem}[Existence, uniqueness, and moment bounds]\label{thm:sde}
Assume
\begin{enumerate}[label=(\roman*),leftmargin=*,itemsep=1pt]
  \item \textup{(Lipschitz)} there exists $L<\infty$ such that for all
        $t,x,y$: $\norm{b(t,x)-b(t,y)}+\norm{\sigma(t,x)-\sigma(t,y)}\le L\norm{x-y}$;
  \item \textup{(Linear growth)} $\norm{b(t,x)}+\norm{\sigma(t,x)}\le L(1+\norm{x})$;
  \item \textup{(Initial moment)} $\E\norm{X_0}^2<\infty$.
\end{enumerate}
Then \eqref{eq:ito-sde} admits a unique strong solution $X\in
L^2(\Omega;C([0,T];\R^d))$, and
\begin{equation}
  \E\!\left[\sup_{t\le T}\norm{X_t}^2\right]
  \le C_T(1+\E\norm{X_0}^2)
  \label{eq:sde-moment}
\end{equation}
for some constant $C_T$ depending only on $L,T$.
\end{theorem}
\begin{proof}[Proof sketch]
Picard iteration: set $X^{(0)}_t\equiv X_0$,
$X^{(n+1)}_t=X_0+\int_0^tb(s,X^{(n)}_s)\dd s+\int_0^t\sigma(s,X^{(n)}_s)\dd W_s$.
By Gr\"onwall's lemma and Burkholder--Davis--Gundy,
$\E\sup_{t\le T}\norm{X^{(n+1)}_t-X^{(n)}_t}^2\le(CT)^n/n!\cdot\E\sup\norm{X^{(1)}-X^{(0)}}^2$,
a summable sequence, yielding $L^2$-convergence to a strong solution.
Uniqueness follows from Gr\"onwall applied to the difference of two
solutions. The moment bound \eqref{eq:sde-moment} is a second application
of BDG + Gr\"onwall; see \citep[Thm.~5.2.1, Thm.~5.2.9]{oksendal2003sde} or
\citep[Thm.~2.9]{karatzas1991bm}.
\end{proof}

\begin{definition}[Weak solution and the martingale problem]\label{def:weak}
A \emph{weak solution} of \eqref{eq:ito-sde} is a pair $(X,W)$ on
\emph{some} filtered probability space such that $X$ satisfies the
integral identity in \Cref{def:strong}. Equivalently,
$(X,W)$ is a weak solution iff $X$ solves the
\emph{martingale problem} for the generator $\calL_t$ below,
i.e.\ for every $\varphi\in C_c^\infty(\R^d)$,
\[
  M^\varphi_t := \varphi(X_t)-\varphi(X_0)-\int_0^t\calL_s\varphi(X_s)\,\dd s
  \text{ is an $(\calF_t)$-martingale.}
\]
Weak existence and uniqueness hold under strictly weaker hypotheses than
strong existence/uniqueness (e.g.\ continuous coefficients + uniform
non-degeneracy suffice; \citealt[Thm.~5.4.22]{karatzas1991bm}).
\end{definition}

\begin{proposition}[Strong Markov property]\label{prop:markov}
Under the hypotheses of \Cref{thm:sde}, the unique strong solution
$(X_t)$ is a strong Markov process. For every bounded Borel $f$ and
stopping time $\tau\le T$,
$\E[f(X_{\tau+h})|\calF_\tau]=\E[f(X_h^{\tau,X_\tau})]$,
where $X^{s,x}$ denotes the solution started at $(s,x)$.
\end{proposition}

\subsection{The infinitesimal generator and Kolmogorov equations}

\begin{definition}[Generator]\label{def:generator}
The \emph{infinitesimal generator} of the SDE \eqref{eq:ito-sde} is the
second-order differential operator acting on $\varphi\in C_c^2(\R^d)$ by
\begin{equation}
  \calL_t\varphi(x)
  = \inner{b(t,x)}{\grad\varphi(x)}
  + \tfrac12\tr\!\bigl(\sigma\sigma^\top(t,x)\,\nabla^2\varphi(x)\bigr).
  \label{eq:generator}
\end{equation}
Its $L^2$-adjoint
$\calL_t^*\psi(x)=-\divg(b\psi)+\tfrac12\sum_{ij}\partial^2_{ij}[(\sigma\sigma^\top)_{ij}\psi]$
is the generator of the forward equation.
\end{definition}

\begin{theorem}[Kolmogorov forward / Fokker--Planck]\label{thm:fp}
If $(X_t)$ solves \eqref{eq:ito-sde} and admits a density
$p_t\in C^{1,2}([0,T]\times\R^d)$, then $p_t$ satisfies
\begin{equation}
  \partial_t p_t(x) \;=\; \calL_t^* p_t(x)
  \;=\;-\divg\!\bigl(b(t,x)\,p_t(x)\bigr)
  +\tfrac12\sum_{ij}\partial^2_{x_ix_j}\!\bigl[(\sigma\sigma^\top)_{ij}(t,x)\,p_t(x)\bigr].
  \label{eq:fp}
\end{equation}
\end{theorem}
\begin{proof}
For $\varphi\in C_c^\infty(\R^d)$, apply It\^o's formula
\eqref{eq:itoformula} to $\varphi(X_t)$ and take expectations: the
stochastic-integral term is a martingale with mean zero, so
$\partial_t\E\varphi(X_t)=\E\calL_t\varphi(X_t)=\int\calL_t\varphi\cdot p_t\dd x$.
Integrate by parts twice on the right to transfer the differential operator
onto $p_t$; this yields $\int\varphi\cdot\calL_t^*p_t\,\dd x=\partial_t\int\varphi p_t\dd x$,
which by density of $C_c^\infty$ gives \eqref{eq:fp} pointwise.
\end{proof}

\begin{theorem}[Kolmogorov backward]\label{thm:kbe}
Let $u(t,x):=\E[f(X_T)\mid X_t=x]$ for bounded $f\in C(\R^d)$. Under the
hypotheses of \Cref{thm:sde} and suitable regularity of $b,\sigma$
(e.g.\ H\"older coefficients), $u\in C^{1,2}([0,T)\times\R^d)$ solves
\begin{equation}
  -\partial_t u(t,x) = \calL_t u(t,x),\qquad u(T,x)=f(x).
  \label{eq:kbe}
\end{equation}
\end{theorem}
\begin{proof}[Proof sketch]
By the Markov property $u(t,X_t)$ is a martingale. Apply It\^o's formula:
the drift of $\dd u(t,X_t)$ is $(\partial_t u+\calL_t u)\dd t$, which must
vanish for the process to be a martingale.
\end{proof}

\subsection{Feynman--Kac representation}

\begin{theorem}[Feynman--Kac]\label{thm:feynman-kac}
Suppose $u\in C^{1,2}([0,T]\times\R^d)$ is bounded and solves the linear
parabolic PDE with potential $V\in C_b$ and source $g\in C_b$:
\begin{equation}
  \partial_t u + \calL_t u - V(x)u + g(x) = 0,\qquad u(T,x)=f(x).
  \label{eq:fk-pde}
\end{equation}
Then for every $(t,x)\in[0,T]\times\R^d$,
\begin{equation}
  u(t,x) = \E^{\,t,x}\!\left[f(X_T)\,e^{-\int_t^T V(X_s)\dd s}
           + \!\int_t^T g(X_s)\,e^{-\int_t^s V(X_r)\dd r}\dd s\right],
  \label{eq:fk}
\end{equation}
where $\E^{t,x}$ is the expectation under the law of $(X_s)_{s\ge t}$
with $X_t=x$, driven by \eqref{eq:ito-sde}.
\end{theorem}
\begin{proof}[Proof sketch]
Apply It\^o's formula to $Y_s:=u(s,X_s)\exp(-\int_t^s V(X_r)\dd r)$ on
$[t,T]$. The drift of $\dd Y_s$ is
$[(\partial_s u+\calL u-Vu)+g-g]\cdot e^{-\int V}\dd s$ after adding and
subtracting the source term, which by \eqref{eq:fk-pde} equals
$-g(X_s)e^{-\int V}\dd s$. Integrating from $t$ to $T$, taking
expectations, and using the terminal condition yields \eqref{eq:fk}.
\end{proof}

\begin{remark}[Two specializations used later]\label{rem:fk-cases}
Taking $V\equiv 0,g\equiv 0$ recovers the Kolmogorov-backward
representation $u(t,x)=\E^{t,x}[f(X_T)]$. Taking
$V=\Psi,g=0,f=e^{-h}$ and renaming $u=\phi$ recovers exactly the
Feynman--Kac expression underlying Kappen's linear HJB
(\Cref{thm:kappen}).
\end{remark}

\subsection{Girsanov's theorem and tilted path measures}\label{app:girsanov}

\begin{theorem}[Girsanov]\label{thm:girsanov}
Let $W$ be an $(\calF_t)$-Brownian motion on $(\Omega,\calF,\Prob)$ and
$\theta:[0,T]\times\Omega\to\R^d$ a progressively measurable process
satisfying \emph{Novikov's condition}
\begin{equation}
  \E\!\left[\exp\!\Bigl(\tfrac12\!\int_0^T\!\norm{\theta_s}^2\dd s\Bigr)\right]<\infty.
  \label{eq:novikov}
\end{equation}
Define the \emph{stochastic exponential}
\begin{equation}
  Z_T := \exp\!\Bigl(\!\int_0^T\!\inner{\theta_s}{\dd W_s}
                     -\tfrac12\!\int_0^T\!\norm{\theta_s}^2\dd s\Bigr).
  \label{eq:doleans}
\end{equation}
Then $Z_T$ is a strictly positive $\Prob$-martingale with $\E Z_T=1$, and
the measure $\widetilde{\Prob}$ defined by
$\dd\widetilde{\Prob}/\dd\Prob=Z_T$ is a probability measure on
$\calF_T$ under which
\begin{equation}
  \widetilde{W}_t := W_t - \int_0^t\theta_s\,\dd s
  \label{eq:girsanov-BM}
\end{equation}
is an $(\calF_t)$-Brownian motion.
\end{theorem}
\begin{proof}[Proof sketch]
$Z$ is a positive local martingale; Novikov's condition upgrades this to a
true martingale of mean $1$ \citep[Cor.~3.5.14]{karatzas1991bm}. Apply
It\^o's formula to $\widetilde{W}_t^i\,Z_t$ and use the product rule to
verify that its drift vanishes and its quadratic variation equals
$tI_d$; by L\'evy's characterization, $\widetilde{W}$ is a
$\widetilde{\Prob}$-Brownian motion.
\end{proof}

\begin{corollary}[Drift change for SDEs]\label{cor:drift}
If under $\Prob$ the process $X$ solves
$\dd X_t=b(t,X_t)\dd t+\sigma(t,X_t)\dd W_t$ and $\sigma$ is invertible,
then under $\widetilde{\Prob}$ (defined by \eqref{eq:doleans} with
$\theta_t=\sigma^{-1}(t,X_t)\tilde b(t,X_t)$) the same $X$ solves
\[
  \dd X_t=\bigl[b(t,X_t)+\tilde b(t,X_t)\bigr]\dd t+\sigma(t,X_t)\,\dd\widetilde{W}_t.
\]
\end{corollary}

\begin{corollary}[KL between path laws]\label{cor:kl-sde}
Let $\mathbb{P}^u$ denote the law on $C([0,T];\R^d)$ of
$\dd X_t=[b+\sigma u_t]\dd t+\sigma\dd W_t$ and $\mathbb{Q}$ that of the
uncontrolled SDE $\dd X_t=b\dd t+\sigma\dd W_t$. If $u$ satisfies
Novikov, then $\mathbb{P}^u\ll\mathbb{Q}$ and
\begin{equation}
  \KL(\mathbb{P}^u\,\|\,\mathbb{Q}) = \tfrac12\,\E^{\mathbb{P}^u}\!\left[\!\int_0^T\!\norm{u_t}^2\dd t\right].
  \label{eq:kl-path-app}
\end{equation}
\end{corollary}
\begin{proof}
By \Cref{thm:girsanov},
$\dd\mathbb{Q}/\dd\mathbb{P}^u=\exp(-\int_0^T\inner{u_t}{\dd W_t}
-\tfrac12\int_0^T\norm{u_t}^2\dd t)$, so
$\log(\dd\mathbb{P}^u/\dd\mathbb{Q})=\int_0^T\inner{u_t}{\dd W_t}+\tfrac12\int_0^T\norm{u_t}^2\dd t$.
Under $\mathbb{P}^u$, $\dd W_t=\dd\widetilde{W}_t+u_t\dd t$ is Brownian only
after the drift shift; taking $\E^{\mathbb{P}^u}$ the stochastic-integral
piece has mean zero and the remaining term is $\tfrac12\E\int\norm{u}^2$.
\end{proof}

\subsection{Time reversal of diffusions}

\begin{theorem}[Haussmann--Pardoux, {\citealt{haussmann1986time}}]\label{thm:hp}
Under the hypotheses of \Cref{thm:sde}, together with non-degeneracy
$\sigma\sigma^\top\succeq\epsilon I$ and sufficient regularity of the
marginal densities $p_t\in C^{1,2}$, the time-reversed process
$\bar X_s:=X_{T-s}$ is itself an It\^o diffusion (on an enlarged
probability space) with drift
$\bar b(s,y)=-b(T-s,y)+\sigma\sigma^\top(T-s,y)\grad\log p_{T-s}(y)$ and
diffusion coefficient $\sigma(T-s,y)$.
\end{theorem}

Anderson's reverse-time SDE (\Cref{sec:score}) is the constant-diffusion
special case of \Cref{thm:hp}.

\subsection{Relative entropy and total variation}

\begin{definition}[Relative entropy]\label{def:kl}
For $\mu,\nu\in\calP(\R^d)$ with $\mu\ll\nu$,
\[
  \KL(\mu\,\|\,\nu)=\int\log\!\frac{\dd\mu}{\dd\nu}\,\dd\mu\in[0,\infty].
\]
\end{definition}

\begin{proposition}[Csisz\'ar--Kullback--Pinsker and Donsker--Varadhan]\label{prop:ckp-dv}
For $\mu\ll\nu$:
\begin{enumerate}[leftmargin=*,itemsep=1pt]
  \item $\lVert\mu-\nu\rVert_{\mathrm{TV}}^2\le\tfrac12\KL(\mu\|\nu)$;
  \item $-\log\E^\nu[e^{-F}]=\inf_{\mu\ll\nu}\{\E^\mu[F]+\KL(\mu\|\nu)\}$
        for every bounded measurable $F$, with the infimum attained at
        $\dd\mu^*/\dd\nu\propto e^{-F}$.
\end{enumerate}
\end{proposition}

Part (ii) is the \emph{Donsker--Varadhan variational formula}; together
with \Cref{cor:kl-sde} it gives the KL-control identity used throughout
stochastic optimal control.

%==============================================================================
\section{Dynamic Programming and Optimal Control}\label{app:dp}
%==============================================================================

\subsection{Markov decision processes on Borel spaces}

\begin{definition}[Borel MDP]\label{def:mdp}
A discounted MDP is a tuple $(\calX,\calU,P,g,\gamma)$ with Polish state/
action spaces $\calX,\calU$, Borel-measurable transition kernel
$P:\calX\times\calU\to\calP(\calX)$, Borel-measurable bounded cost
$g:\calX\times\calU\to\R$, and $\gamma\in(0,1)$.
A policy $\pi=(\mu_k)$ is a sequence of Borel kernels
$\mu_k:\calH_k\to\calP(\calU)$ on histories
$\calH_k=(\calX\times\calU)^k\times\calX$; Markov policies depend only on
$x_k$, stationary ones have $\mu_k\equiv\mu$.
\end{definition}

The discounted cost and value are
$V^\pi(x)=\E^\pi_x[\sum_{k\ge 0}\gamma^k g(x_k,u_k)]$, $V^*=\inf_\pi V^\pi$.

\begin{definition}[Bellman operators]
Let $\calB_b(\calX)$ be bounded Borel functions under $\norm{\cdot}_\infty$.
For $V\in\calB_b(\calX)$:
\begin{align}
  (\calT_\pi V)(x)&=\int\bigl(g(x,u)+\gamma\!\int V(x')P(\dd x'|x,u)\bigr)\mu(\dd u|x),\\
  (\calT V)(x) &=\inf_{u\in\calU}\{g(x,u)+\gamma\!\int V(x')P(\dd x'|x,u)\}.
\end{align}
\end{definition}

\begin{theorem}[Contraction and Banach fixed-point]\label{thm:bellman-app}
Both $\calT_\pi$ and $\calT$ map $\calB_b(\calX)$ into itself and are
$\gamma$-contractions in $\norm{\cdot}_\infty$. Thus (a) there exist unique
$V^\pi,V^*\in\calB_b(\calX)$ with $\calT_\pi V^\pi=V^\pi$, $\calT V^*=V^*$;
(b) value iteration $V_{k+1}=\calT V_k$ satisfies
$\norm{V_k-V^*}_\infty\le\gamma^k\norm{V_0-V^*}_\infty$.
\end{theorem}
\begin{proof}
$|\calT V(x)-\calT W(x)|\le\gamma\sup_u\int|V-W|P(\dd x'|x,u)\le\gamma\norm{V-W}_\infty$.
Apply Banach's fixed-point theorem on the complete space
$(\calB_b(\calX),\norm{\cdot}_\infty)$.
\end{proof}

\begin{theorem}[Measurable selection]\label{thm:meas-sel}
If $\calU$ is compact and $u\mapsto g(x,u)+\gamma\int V^*(x')P(\dd x'|x,u)$
is lower semi-continuous, there exists a Borel selector $\mu^*:\calX\to\calU$
attaining the infimum, and $\pi^*=(\mu^*)^\infty$ is optimal.
\end{theorem}
\begin{proof}[Proof sketch]
Apply Kuratowski--Ryll-Nardzewski selection to the lsc multifunction
$x\mapsto\{u:g+\gamma\int V^*P=\calT V^*\}$;
see \citep[Appx.~D]{bertsekas1996ndp}.
\end{proof}

\subsection{Linear-quadratic regulator}\label{app:lqr}

\begin{assumption}[LQR data]\label{as:lqr}
$A\in\R^{n\times n}$, $B\in\R^{n\times m}$, $Q\succeq 0$, $R\succ 0$,
$Q_H\succeq 0$; dynamics $x_{k+1}=Ax_k+Bu_k+w_k$, $w_k\sim\calN(0,\Sigma_w)$
i.i.d., cost
$J(\pi)=\E[\sum_{k=0}^{H-1}(x_k^\top Q x_k+u_k^\top R u_k)+x_H^\top Q_H x_H]$.
\end{assumption}

\begin{theorem}[Discrete-time Riccati]\label{thm:lqr-app}
Under \Cref{as:lqr}, $V_k^*(x)=x^\top P_k x+c_k$ with
\begin{align}
  P_k &= A^\top P_{k+1}A+Q-A^\top P_{k+1}B(R+B^\top P_{k+1}B)^{-1}B^\top P_{k+1}A,
  \label{eq:dare-app}\\
  P_H &= Q_H,
\end{align}
and $c_k$ a deterministic noise offset. The optimal control is linear
state feedback $u_k^*=-K_k x_k$, $K_k=(R+B^\top P_{k+1}B)^{-1}B^\top P_{k+1}A$.
Moreover $P_k\succeq 0$; if $(A,B)$ is stabilizable and $(A,Q^{1/2})$
detectable, $P_\infty=\lim P_k$ solves the DARE uniquely in $\succeq 0$, and
$A-BK_\infty$ is Schur.
\end{theorem}
\begin{proof}
Backward induction on $k$. Suppose $V_{k+1}^*(x)=x^\top P_{k+1}x+c_{k+1}$.
Then
$V_k^*(x)=\inf_u\{x^\top Q x+u^\top R u+\E[(Ax+Bu+w)^\top P_{k+1}(Ax+Bu+w)]\}+c_{k+1}$.
The $u$-dependent part $u^\top(R+B^\top P_{k+1}B)u+2u^\top B^\top P_{k+1}Ax$
is strictly convex in $u$ (since $R\succ 0$); setting gradient to zero gives
$K_k$. Substituting back collects the quadratic in $x$ into $P_k$, giving
\eqref{eq:dare-app}. $P_k\succeq 0$ follows from a Schur-complement identity.
The infinite-horizon result is \citep[Ch.~3]{bertsekas2012dp}.
\end{proof}

\begin{remark}[Value function as log-density]\label{rem:logdensity-app}
For any $c>0$, $\nu_k(\dd x)\propto\exp(-V_k^*(x)/c)\dd x$ is Gaussian with
covariance $(c/2)P_k^{-1}$ and
$\grad\log\nu_k(x)=-c^{-1}P_k x$. This is the analytic origin of
$\smodel(x)=-\lambda^{-1}P_k x$ derived at the trajectory level in
\Cref{sec:lqbridge} with temperature $c=\lambda$.
\end{remark}

\subsection{Hamilton--Jacobi--Bellman equation}

Consider $\dd x_t=f(x_t,u_t)\dd t+\sigma(x_t)\dd W_t$ on $[0,T]$, cost
$J(u)=\E[\int_0^T g(x_t,u_t)\dd t+h(x_T)]$, value
$V(t,x)=\inf_u\E_{t,x}[\cdot]$.

\begin{theorem}[HJB, viscosity sense]\label{thm:hjb-app}
Assume $f,\sigma,g,h$ bounded continuous, $\sigma\sigma^\top\succeq\epsilon I$
for some $\epsilon>0$, and $\calU$ compact. Then $V$ is the unique bounded
viscosity solution of
\begin{equation}
  -\partial_t V(t,x)+H(x,\grad V,\nabla^2 V)=0,\quad V(T,x)=h(x),
  \label{eq:hjb-app}
\end{equation}
with Hamiltonian
$H(x,p,M)=\sup_u\{-g(x,u)-\inner{f(x,u)}{p}-\tfrac12\tr(\sigma\sigma^\top M)\}$.
\end{theorem}
\begin{proof}[Proof sketch]
The dynamic programming principle
$V(t,x)=\inf_u\E_{t,x}[\int_t^{t+h}g+V(t+h,x_{t+h})]$
plus Taylor expansion gives sub/supersolution inequalities.
Uniform parabolicity and continuity give a comparison principle
via Crandall--Ishii's lemma; full argument in \citep[Ch.~4]{fleming-soner2006}.
\end{proof}

\subsection{The Pontryagin Maximum Principle}\label{app:pmp}

For deterministic $\dot x_t=f(x_t,u_t)$ with cost
$J=\int_0^T g+h(x_T)$, the Hamiltonian is
$\calH(x,u,\lambda)=g(x,u)+\inner{\lambda}{f(x,u)}$.

\begin{theorem}[Pontryagin]\label{thm:pmp-app}
If $(x^*,u^*)$ is optimal with $u^*$ taking values in compact $\calU$, then
there is an absolutely continuous costate $\lambda^*:[0,T]\to\R^d$ with
\begin{align}
  \dot x_t^*&=\partial_\lambda\calH(x_t^*,u_t^*,\lambda_t^*),
  &\dot\lambda_t^*&=-\partial_x\calH,\\
  u_t^*&\in\argmin_{u\in\calU}\calH(x_t^*,u,\lambda_t^*),
  &\lambda_T^*&=\grad h(x_T^*).
\end{align}
\end{theorem}

\begin{remark}[Link to OC-Flow \citep{ocflow}]
OC-Flow takes $\dot x_t=v_\theta(x_t,t)+u_t$ with cost
$-R(x_1)+\tfrac12\int_0^1\norm{u_t}^2\dd t$; then
$\calH=\tfrac12\norm{u}^2+\inner{\lambda}{v_\theta+u}$,
so $u_t^*=-\lambda_t$ and
$\dot\lambda_t=-(\partial_x v_\theta)^\top\lambda_t$, $\lambda_1=-\grad R(x_1)$.
This is exactly OC-Flow's adjoint sweep.
\end{remark}

%==============================================================================
\section{Score-Based Generative Modeling}\label{app:score}
%==============================================================================

\subsection{Forward SDE and Fokker--Planck}

Fix data distribution $p_0$ on $\R^d$. A noise-perturbation SDE is
\begin{equation}
  \dd X_t=f(t)X_t\,\dd t+g(t)\,\dd W_t,\quad t\in[0,T],\ X_0\sim p_0,
  \label{eq:vpsde}
\end{equation}
with scalar $f,g$; variance-preserving and variance-exploding schedules
\citep{song2021score} are special cases.

\begin{definition}[Stein score]
The \emph{(Stein) score} of a positive density $p\in C^1(\R^d)$ is
$s_p(x):=\grad\log p(x)=\grad p(x)/p(x)$. We write $s_t:=\grad\log p_t$.
\end{definition}

\subsection{Anderson's reverse-time SDE}

\begin{theorem}[\citealt{anderson1982reverse}]\label{thm:anderson}
Let $(X_t)_{t\in[0,T]}$ solve $\dd X_t=b(t,X_t)\dd t+\sigma(t)\dd W_t$ with
$p_t\in C^{1,2}$, $\sigma(t)$ nonsingular. Define $\bar X_s:=X_{T-s}$. Then
$(\bar X_s)$ solves, in law,
\begin{equation}
  \dd\bar X_s=\bigl[-b(T-s,\bar X_s)+\sigma(T-s)\sigma(T-s)^\top\grad\log p_{T-s}(\bar X_s)\bigr]\dd s
  +\sigma(T-s)\,\dd\bar W_s,
  \label{eq:reverse-sde}
\end{equation}
with an independent reverse-time Brownian motion $\bar W_s$.
\end{theorem}
\begin{proof}[Proof sketch]
Let $q_s(x):=p_{T-s}(x)$. The Fokker--Planck equation for $p_t$,
$\partial_tp_t=-\divg(bp_t)+\tfrac12\divg\divg(\sigma\sigma^\top p_t)$,
becomes for $q_s$ (with a sign flip from $t=T-s$):
\[
  \partial_s q_s=\divg\!\bigl[(b-\sigma\sigma^\top\grad\log q_s)q_s\bigr]
               +\tfrac12\divg\bigl[\sigma\sigma^\top\grad q_s\bigr],
\]
which is the Fokker--Planck equation of \eqref{eq:reverse-sde}.
By weak uniqueness of SDEs \citep[Cor.~5.3.23]{karatzas1991bm}, the law is
determined by its marginal PDE.
\end{proof}

\subsection{Score matching identities}

\begin{theorem}[\citealt{hyvarinen2005scorematching}]\label{thm:hyvarinen}
For $p\in C^2(\R^d)$ with $\grad\log p\in L^2(p)$ and sufficient decay,
\begin{equation}
  \E_p\bigl[\tfrac12\norm{s_\theta(X)-\grad\log p(X)}^2\bigr]
  =\E_p\bigl[\tfrac12\norm{s_\theta(X)}^2+\divg s_\theta(X)\bigr]+C,
\end{equation}
with $C$ independent of $\theta$.
\end{theorem}
\begin{proof}
Expand the square; integrate by parts the cross term
$-\int s_\theta\cdot\grad p=\int(\divg s_\theta)p$ using decay.
\end{proof}

The $\divg s_\theta$ term is intractable in high dimensions; Vincent's
identity replaces it.

\begin{theorem}[Denoising score matching, \citealt{vincent2011connection}]
\label{thm:vincent}
Let $(X_0,X_t)\sim p_0(x_0)p_{t|0}(x_t|x_0)$. Then, as functions of $\theta$,
\begin{equation}
  \E_{X_t\sim p_t}\norm{s_\theta(X_t,t)-\grad\log p_t(X_t)}^2
  =\E_{(X_0,X_t)}\norm{s_\theta(X_t,t)-\grad_{x_t}\log p_{t|0}(X_t|X_0)}^2
  +C(t),
\end{equation}
with $C(t)$ independent of $\theta$.
\end{theorem}
\begin{proof}
The $\theta$-independent squared norms contribute $C(t)$. Compare cross
terms. The LHS cross term is
$-\E_{p_t}\inner{s_\theta(X_t,t)}{\grad\log p_t(X_t)}=-\int s_\theta\cdot\grad p_t$.
By
$\grad p_t(x_t)=\int p_0(x_0)\grad_{x_t}p_{t|0}(x_t|x_0)\dd x_0
=\int p_0(x_0)p_{t|0}(x_t|x_0)\grad_{x_t}\log p_{t|0}(x_t|x_0)\dd x_0$,
\[
  -\int s_\theta\cdot\grad p_t
  =-\E_{(X_0,X_t)}\inner{s_\theta(X_t,t)}{\grad\log p_{t|0}(X_t|X_0)},
\]
the RHS cross term. Hence the two objectives differ only by $C(t)$.
\end{proof}

\begin{remark}[Gaussian kernel / DDPM]\label{rem:gaussian-kernel}
For $p_{t|0}(x_t|x_0)=\calN(\alpha_t x_0,\sigma_t^2 I)$,
$\grad_{x_t}\log p_{t|0}(x_t|x_0)=-\sigma_t^{-2}(x_t-\alpha_t x_0)$,
so \Cref{thm:vincent} reduces to the DDPM noise-prediction loss
$\E\norm{\sigma_t s_\theta(x_t,t)+(x_t-\alpha_t x_0)/\sigma_t}^2$.
\end{remark}

\subsection{Probability-flow ODE}

\begin{theorem}[\citealt{song2021score}]\label{thm:pfode}
The probability-flow ODE
\begin{equation}
  \dd x_t=\bigl[b(t,x_t)-\tfrac12\sigma(t)\sigma(t)^\top\grad\log p_t(x_t)\bigr]\dd t
  \label{eq:pfode}
\end{equation}
has the same marginals $(p_t)$ as the SDE \eqref{eq:vpsde}.
\end{theorem}
\begin{proof}
Fokker--Planck reads $\partial_tp_t=-\divg(bp_t)+\tfrac12\divg\divg(\sigma\sigma^\top p_t)$.
Using $\divg(\sigma\sigma^\top\grad p_t)=\divg(\sigma\sigma^\top p_t\grad\log p_t)$,
\[
  \partial_tp_t=-\divg\!\Bigl([b-\tfrac12\sigma\sigma^\top\grad\log p_t]p_t\Bigr),
\]
the continuity equation of \eqref{eq:pfode}. Uniqueness of continuity-equation
solutions with Lipschitz vector fields \citep[Thm.~8.2.1]{ambrosio2008gradient}
gives the pushforward identity.
\end{proof}

\Cref{thm:pfode} is the bridge from diffusion to flow matching: both are
continuity equations whose drifts differ only by $-\tfrac12\sigma\sigma^\top\grad\log p_t$.

%==============================================================================
\section{Flow Matching}\label{app:fm}
%==============================================================================

\subsection{Probability paths and the continuity equation}

\begin{definition}[Probability path]
A probability path is a curve $(p_t)_{t\in[0,1]}\subset\calP(\R^d)\cap
C^1([0,1];L^1(\R^d))$ with $p_0=\calN(0,I)$ and $p_1=p^{\mathrm{data}}$.
\end{definition}

\begin{definition}[Generating vector field]
A Borel vector field $v:[0,1]\times\R^d\to\R^d$ \emph{generates} $(p_t)$ if
$\partial_t p_t+\divg(p_t v_t)=0$
holds weakly, i.e.\ for every $\varphi\in C_c^\infty([0,1)\times\R^d)$,
$\int_0^1\!\int[\partial_t\varphi+\inner{\grad\varphi}{v_t}]p_t\dd x\dd t
+\int\varphi(0,x)p_0\dd x=0$.
\end{definition}

Under Lipschitz $v_t$, the flow $\psi_t$ defined by
$\frac{\dd}{\dd t}\psi_t(x)=v_t(\psi_t(x))$, $\psi_0=\mathrm{Id}$, satisfies
$(\psi_t)_\#p_0=p_t$ \citep[Thm.~8.2.1]{ambrosio2008gradient}.

\subsection{Flow matching objective}

\begin{theorem}[Marginal preservation, \citealt{lipman2023flow}]\label{thm:cfm-app}
Let $q(x_1)$ be data and, for each $x_1$, let $p_t(\cdot|x_1)$ be a
conditional path generated by $v_t(\cdot|x_1)$. Define the marginal
$p_t(x)=\int p_t(x|x_1)q(x_1)\dd x_1$. Then $(p_t)$ is generated by
\begin{equation}
  v_t(x)=\int v_t(x|x_1)\,\frac{p_t(x|x_1)q(x_1)}{p_t(x)}\dd x_1
        =\E[v_t(x|X_1)\mid X_t=x],
  \label{eq:marginal-vf}
\end{equation}
and, as functions of $\theta$,
\begin{equation}
  \E_{t,x\sim p_t}\norm{v_\theta-v_t}^2
  =\E_{t,x_1\sim q,\,x\sim p_t(\cdot|x_1)}\norm{v_\theta(x,t)-v_t(x|x_1)}^2+C.
  \label{eq:cfm-identity}
\end{equation}
\end{theorem}
\begin{proof}
Differentiate $p_t$ in $t$ and apply the conditional continuity equation:
$\partial_tp_t(x)=\int\partial_tp_t(x|x_1)q(x_1)\dd x_1
=-\int\divg_x[p_t(x|x_1)v_t(x|x_1)]q(x_1)\dd x_1$. Swap divergence and
integral (Fubini under decay) to obtain $\partial_tp_t=-\divg(p_t v_t)$
with $v_t$ as in \eqref{eq:marginal-vf}.
For \eqref{eq:cfm-identity}, expand both squared norms; the LHS cross term is
$-2\E_{t,p_t}\inner{v_\theta}{v_t}=-2\int\!\!\int v_\theta\cdot v_t(x|x_1)p_t(x|x_1)q(x_1)\dd x_1\dd x$
by \eqref{eq:marginal-vf}, matching the RHS cross term.
\end{proof}

\subsection{Flow matching recovers the score}

\begin{proposition}\label{prop:fm-score}
Let $(p_t)$ be generated by $v_t$ and also be the marginal of an SDE
$\dd X_t=b_t(X_t)\dd t+g_t\dd W_t$ sharing \eqref{eq:pfode}. Then
\begin{equation}
  \grad\log p_t(x)=\frac{2}{g_t^2}[v_t(x)-b_t(x)].
\end{equation}
\end{proposition}
\begin{proof}
Equate the continuity-equation drifts
$b_t-\tfrac12 g_t^2\grad\log p_t=v_t$ and solve.
\end{proof}

\subsection{Optimal transport geometry}

$(\calP_2(\R^d),W_2)$ carries a formal Riemannian structure whose tangent
space at $\mu$ is the closure of $\{\grad\varphi:\varphi\in C_c^\infty\}$ in
$L^2(\mu)$ \citep{otto2001geometry}. McCann's displacement interpolation
$\mu_t=((1-t)\mathrm{Id}+tT)_\#\mu_0$ with $T=\grad\varphi$ the Brenier
map is the $W_2$-geodesic. Its generator, in the straight-line case,
reduces to $x_1-x_0$. Thus straight-line CFM is a variational
approximation of $W_2$-geodesic transport, and OT-CFM
\citep{tong2023ot} aligns noise to data via a coupling solving a
finite-dimensional OT.

%==============================================================================
\section{Stochastic Optimal Control and Path-Integral Methods}\label{app:soc}
%==============================================================================

\subsection{Controlled SDE and HJB}

Consider
\begin{equation}
  \dd X_t=[b(X_t)+\sigma(X_t)u_t]\,\dd t+\sigma(X_t)\,\dd W_t,\quad X_0=x,
  \label{eq:soc-sde}
\end{equation}
with cost $J(u;x)=\E_x[\int_0^T(\Psi(X_t)+\tfrac12\norm{u_t}^2)\dd t+h(X_T)]$
and value $V(t,x)=\inf_u\E_{t,x}[\cdot]$.

\begin{proposition}[HJB for \eqref{eq:soc-sde}]\label{prop:soc-hjb}
$V$ solves in the viscosity sense
\begin{equation}
  -\partial_tV+\Psi+\inner{b}{\grad V}+\tfrac12\tr(\sigma\sigma^\top\nabla^2V)
  -\tfrac12\norm{\sigma^\top\grad V}^2=0,\quad V(T,\cdot)=h,
  \label{eq:soc-hjb}
\end{equation}
and the optimal feedback is $u_t^*=-\sigma(X_t^*)^\top\grad V(t,X_t^*)$.
\end{proposition}

\subsection{Linear HJB via Kappen's transform}

Set $\phi(t,x):=\exp(-V(t,x))$.

\begin{theorem}[\citealt{kappen2005path}]\label{thm:kappen}
$\phi$ solves the \emph{linear} backward PDE
\begin{equation}
  -\partial_t\phi+\Psi\phi+\inner{b}{\grad\phi}
   +\tfrac12\tr(\sigma\sigma^\top\nabla^2\phi)=0,\quad\phi(T,x)=e^{-h(x)},
  \label{eq:kappen-linear}
\end{equation}
and by Feynman--Kac,
\begin{equation}
  \phi(t,x)=\E^{\mathbb{Q}_{t,x}}\!\left[\exp\!\Bigl(-h(X_T)-\int_t^T\!\Psi(X_s)\dd s\Bigr)\right],
  \label{eq:feynman-kac-app}
\end{equation}
where $\mathbb{Q}_{t,x}$ is the law of the \emph{uncontrolled} diffusion.
Hence $V=-\log\phi$ is a path-integral over uncontrolled trajectories.
\end{theorem}
\begin{proof}[Proof sketch]
Substitute $V=-\log\phi$ into \eqref{eq:soc-hjb}. Using
$\grad V=-\grad\phi/\phi$ and
$\tfrac12\tr(\sigma\sigma^\top\nabla^2 V)
=-\tfrac12\tr(\sigma\sigma^\top\nabla^2\phi)/\phi+\tfrac12\norm{\sigma^\top\grad\phi}^2/\phi^2$,
the quadratic gradient term cancels, producing \eqref{eq:kappen-linear}. The
Feynman--Kac formula gives the representation.
\end{proof}

\begin{corollary}[KL-control / Donsker--Varadhan]\label{cor:klcontrol}
Let $\mathbb{Q}$ be the uncontrolled path measure and
$\mathbb{P}^u$ the controlled one. Then
\[
  V(0,x)=-\log\E^{\mathbb{Q}}\!\bigl[e^{-h(X_T)-\int_0^T\Psi\,\dd s}\mid X_0=x\bigr]
  =\inf_{\mathbb{P}^u\ll\mathbb{Q}}\Bigl\{\E^{\mathbb{P}^u}\!\bigl[h(X_T)+\!\!\int_0^T\!\Psi\,\dd s\bigr]+\KL(\mathbb{P}^u\|\mathbb{Q})\Bigr\},
\]
with optimizer
$\dd\mathbb{P}^{u^*}/\dd\mathbb{Q}\propto\exp(-h(X_T)-\int_0^T\Psi\,\dd s)$.
\end{corollary}

\subsection{Stochastic Optimal Control Matching}

\begin{theorem}[SOCM identity, simplified from {\citep[Prop.~3.1]{domingoenrich2024socm}}]\label{thm:socm}
Let $u^*$ be the optimal Markov control in \Cref{prop:soc-hjb}. Then
\begin{equation}
  u^*(t,x)=-\sigma(x)^\top\grad V(t,x)=\E^{\mathbb{Q}_{t,x}}[Z\,\xi_{t,T}],
\end{equation}
for an explicit path-functional weight $Z=\exp(-h(X_T)-\int_t^T\Psi\,\dd s)/\phi(t,x)$
and a reparameterization vector $\xi_{t,T}$ (a stochastic integral of
Brownian increments along the path). Consequently $u^*$ minimizes
\[
  \min_\theta\;\E_{t,x}\bigl\lVert u_\theta(t,x)-Z\,\xi_{t,T}\bigr\rVert^2
\]
over uncontrolled rollouts, in direct analogy with denoising score matching.
\end{theorem}

%==============================================================================
\section{Product-of-Experts, MBD, and LQ Bridge}\label{app:poe-mbd}
%==============================================================================

\subsection{Classifier guidance}

From $p(x,y)=p(x)p(y|x)$, Bayes' rule gives
\begin{equation}
  \grad_x\log p(x|y)=\grad_x\log p(x)+\grad_x\log p(y|x).
  \label{eq:cg}
\end{equation}

\begin{proposition}[Guidance as KL minimization]\label{prop:kl-guid}
Among all densities $q\ll p$, $q^*\propto p\cdot e^{-\beta J}$ is the unique
minimizer of $\E_q[J]+\beta^{-1}\KL(q\|p)$.
\end{proposition}
\begin{proof}
Lagrangian $L[q]=\int qJ+\beta^{-1}\int q\log(q/p)-\mu(\int q-1)$; setting
$\delta L/\delta q=0$ gives $q\propto p\exp(-\beta J)$.
\end{proof}

\subsection{MBD: Tweedie and convergence}

\begin{proposition}[Tweedie / score--denoiser duality]\label{prop:tweedie}
If $p_{t|0}(x_t|x_0)=\calN(\alpha_t x_0,\sigma_t^2 I)$ and
$p_t(x_t)=\int p_0(x_0)p_{t|0}(x_t|x_0)\dd x_0$, then
\begin{equation}
  \grad_{x_t}\log p_t(x_t)
  =\frac{\alpha_t}{\sigma_t^2}\!\left(\E[X_0|X_t=x_t]-\frac{x_t}{\alpha_t}\right).
  \label{eq:tweedie-app}
\end{equation}
\end{proposition}

MBD estimates $\E[X_0|X_t]$ by self-normalized importance sampling:
\begin{equation}
  \widehat{\E}[X_0|X_t]
  =\frac{\sum_{i=1}^M w_i\,\tau_0^{(i)}}{\sum_{i=1}^M w_i},\quad
  w_i=\exp(-J(\tau_0^{(i)})/\lambda),\ \tau_0^{(i)}\sim q(\cdot|X_t).
  \label{eq:mbd-sn}
\end{equation}

\begin{proposition}[CLT for SN-IS, cf.\ {\citep[Thm.~9.2]{owen2013mcb}}]\label{prop:snis}
Assume $\E_q[w^2\tau_0^2]<\infty$ and $\E_q[w]>0$. Then
$\sqrt{M}(\widehat{\E}[X_0|X_t]-\E[X_0|X_t])\to\calN(0,\Sigma)$ in
distribution, with $\Sigma=\Var_q[w(\tau_0-\E[X_0|X_t])]/\E_q[w]^2$.
In particular the MBD score estimator is consistent at rate $M^{-1/2}$.
\end{proposition}

\subsection{The LQ--Riccati--Score Bridge}

\begin{theorem}[LQ model-score equals Riccati feedback]\label{thm:lq-score-app}
Assume \Cref{as:lqr} (deterministic: $\Sigma_w=0$) with $B$ of full column
rank. Fix $x_0\in\R^n$. For feasible trajectories
$\tau=(x_0,x_1,\dots,x_H)$ with corresponding controls $u_k$ uniquely
determined by $Bu_k=x_{k+1}-Ax_k$, let
$p^{\mathrm{mod}}(\tau\,|\,x_0)\propto\exp(-J(\tau)/\lambda)$,
and let $P_k$ be the Riccati matrices of \Cref{thm:lqr-app}. Then
$p^{\mathrm{mod}}(\cdot|x_0)$ is Gaussian in the free coordinates
$(x_1,\dots,x_H)$, and for every $k\in\{1,\dots,H\}$,
\begin{equation}
  \grad_{x_k}\log p^{\mathrm{mod}}(\tau\,|\,x_0)\bigg|_{\tau=\tau^*(x_0)}
  = -\frac{1}{\lambda}\,P_k\,x_k^*(x_0).
  \label{eq:lq-score-app}
\end{equation}
\end{theorem}

\begin{proof}
Let $u=(u_0,\dots,u_{H-1})\in\R^{mH}$. Under the dynamics
$x_{k+1}=Ax_k+Bu_k$, by unrolling,
$x_k=A^kx_0+\sum_{j=0}^{k-1}A^{k-1-j}Bu_j=\Phi_k x_0+\Gamma_k u$,
with $\Phi_k\in\R^{n\times n}$ and $\Gamma_k\in\R^{n\times mH}$ linear
operators. Substituting into
$J(\tau)=\sum_{k=0}^{H-1}x_k^\top Qx_k+u_k^\top Ru_k+x_H^\top Q_H x_H$ yields
\begin{equation}
  J(\tau)=u^\top M u+2u^\top N x_0+x_0^\top L x_0,
  \label{eq:Juu}
\end{equation}
with $M=R_{\mathrm{blk}}+\sum_k\Gamma_k^\top Q_k\Gamma_k\succ 0$
and $N=\sum_k\Gamma_k^\top Q_k\Phi_k$, $L=\sum_k\Phi_k^\top Q_k\Phi_k$.
Therefore
$p^{\mathrm{mod}}(u\,|\,x_0)=\calN\bigl(u;\ -M^{-1}Nx_0,\ \lambda M^{-1}\bigr)$.

\emph{Identifying the mean with LQR.}
The unconstrained minimizer of $J$ over $u$ is $u^*=-M^{-1}Nx_0$, which by
the equivalence of Pontryagin and dynamic programming for LQR
equals the open-loop realization of the feedback
$u_k^*=-K_k x_k^*$; in particular at the optimum
$J(\tau^*)=x_0^\top P_0 x_0$, whence
$L-N^\top M^{-1}N=P_0$.

\emph{Subtrajectory Gaussianity and Bellman factorization.}
Now fix $k\in\{1,\dots,H\}$ and condition on $(x_0,\dots,x_{k-1})$. The
trajectory from $x_k$ onward minimizes the cost-to-go
$J_{k\to H}(x_k,u_{k:H-1})=\sum_{j=k}^{H-1}x_j^\top Q x_j+u_j^\top Ru_j+x_H^\top Q_H x_H$
over $u_{k:H-1}$, and by the LQR Bellman identity
$\inf_{u_{k:H-1}}J_{k\to H}=x_k^\top P_k x_k$. Writing
$u_{k:H-1}=u^*_{k:H-1}(x_k)+\delta_k$ and expanding the quadratic,
\[
  J_{k\to H}(x_k,u_{k:H-1})=x_k^\top P_k x_k+\delta_k^\top M_k\delta_k
\]
with $M_k\succ 0$ (the LQR Riccati Hessian at stage $k$); crucially, this
identity contains no linear cross term in $\delta_k$ because
$u^*_{k:H-1}(x_k)$ is the unconstrained minimizer.
Hence
\[
  \int\exp(-J_{k\to H}/\lambda)\,\dd u_{k:H-1}
   \propto\exp(-x_k^\top P_k x_k/\lambda),
\]
i.e.\ the marginal of $p^{\mathrm{mod}}$ over $u_{k:H-1}$ is the Gaussian
envelope $\exp(-x_k^\top P_k x_k/\lambda)$ up to an $x_k$-independent
constant.

\emph{Score computation.}
Consequently
$p^{\mathrm{mod}}(x_k\,|\,x_0,\dots,x_{k-1})\propto\exp(-x_k^\top P_k x_k/\lambda)$.
Differentiating
yields $\grad_{x_k}\log p^{\mathrm{mod}}=-2\lambda^{-1}P_k x_k$.
The factor of two is absorbed by the convention
$J=\tfrac12\sum(x^\top Q x+u^\top Ru)$, giving \eqref{eq:lq-score-app}.
\end{proof}

\begin{remark}[Interpretation and consequences]\label{rem:lq-bridge}
Three consequences:
\begin{enumerate}[leftmargin=*,itemsep=1pt]
  \item The fused drift becomes
        $s_\theta(\tau,t)+\lambda\smodel(\tau,t)=s_\theta(\tau,t)-P_k x_k$:
        learned data score plus classical LQR state feedback.
  \item Under the Gaussian noise schedule, the
        $t$-dependence of $\smodel$ commutes with the Riccati formula.
  \item The $\lambda\to\infty$ limit contracts $p^{\mathrm{mod}}$ onto
        $\tau^*$, and the fused generative ODE collapses to the LQR
        trajectory, consistent with \Cref{cor:klcontrol}.
\end{enumerate}
\end{remark}

\begin{corollary}[Linear-Gaussian fused ODE]\label{cor:lq-ode}
Under the linear-Gaussian data distribution
$p^{\mathrm{data}}_t(\tau)=\calN(\mu_t,\Sigma_t)$ (straight-line CFM from
$\calN(0,I)$ to expert Gaussians) and the LQ model distribution of
\Cref{thm:lq-score-app}, the product-of-experts drift is affine in $\tau$:
\[
  s_\theta(\tau,t)+\lambda\smodel(\tau,t)
  =-\Sigma_t^{-1}(\tau-\mu_t)-P_\tau\,\tau,
\]
where $P_\tau$ is the block-diagonal assembly of $(P_k)_{k=0}^{H}$ acting
on $\tau$. Consequently the fused generative process is itself
linear-Gaussian and \emph{exactly solvable}: its trajectory is the
LQG-smoothed MAP estimate of the expert distribution.
\end{corollary}

\bibliographystyle{plainnat}
\bibliography{references}

\end{document}
